% =====================================================================
%  TÜBİTAK 2209-A BAŞVURU KILAVUZU — İdris Avlar
%  Derleme (proje kökünden):
%     python kod/kilavuz_ek_uret.py      (Ek A'yı form metninden üretir)
%     pdflatex kilavuz.tex ; pdflatex kilavuz.tex   (iki kez: içindekiler için)
% =====================================================================
\documentclass[11pt,a4paper]{article}
\usepackage[utf8]{inputenc}
\usepackage[T1]{fontenc}
\usepackage{lmodern}
\usepackage[shorthands=off,turkish]{babel}
\usepackage[margin=2.2cm]{geometry}
\usepackage{amsmath,amssymb}
\usepackage{booktabs,longtable,tabularx,array}
\usepackage{enumitem}
\usepackage[table]{xcolor}
\usepackage[most]{tcolorbox}
\usepackage{fancyhdr}
\usepackage{xurl}
\usepackage[hidelinks]{hyperref}
\DeclareUnicodeCharacter{2032}{\ensuremath{'}}
\DeclareUnicodeCharacter{03C3}{\ensuremath{\sigma}}
\DeclareUnicodeCharacter{2212}{\ensuremath{-}}
\DeclareUnicodeCharacter{0394}{\ensuremath{\Delta}}
\DeclareUnicodeCharacter{00B7}{\ensuremath{\cdot}}
\DeclareUnicodeCharacter{00B1}{\ensuremath{\pm}}
\DeclareUnicodeCharacter{00B2}{\ensuremath{^2}}
\DeclareUnicodeCharacter{2075}{\ensuremath{^5}}
\DeclareUnicodeCharacter{03C7}{\ensuremath{\chi}}
\DeclareUnicodeCharacter{03C0}{\ensuremath{\pi}}
\DeclareUnicodeCharacter{00BD}{\ensuremath{\tfrac{1}{2}}}
\DeclareUnicodeCharacter{2265}{\ensuremath{\geq}}
\DeclareUnicodeCharacter{2264}{\ensuremath{\leq}}
\DeclareUnicodeCharacter{2248}{\ensuremath{\approx}}
\DeclareUnicodeCharacter{00D7}{\ensuremath{\times}}
\DeclareUnicodeCharacter{2192}{\ensuremath{\rightarrow}}
\DeclareUnicodeCharacter{03B3}{\ensuremath{\gamma}}
\DeclareUnicodeCharacter{221A}{\ensuremath{\surd}}
\DeclareUnicodeCharacter{03C9}{\ensuremath{\omega}}

\setlength{\parskip}{0.45em}
\setlength{\parindent}{0pt}
\renewcommand{\arraystretch}{1.25}
\pagestyle{fancy}\fancyhf{}
\lhead{\small 2209-A Başvuru Kılavuzu}\rhead{\small İdris Avlar}\cfoot{\thepage}
\setlist{itemsep=2pt,topsep=3pt}
\sloppy

\newtcolorbox{soyle}[1][]{colback=green!5,colframe=green!45!black,title={#1},fonttitle=\bfseries,breakable}
\newtcolorbox{dikkat}[1][]{colback=red!4,colframe=red!60!black,title={#1},fonttitle=\bfseries,breakable}
\newtcolorbox{bilgi}[1][]{colback=blue!4,colframe=blue!50!black,title={#1},fonttitle=\bfseries,breakable}

% Soru-cevap: \S{soru}{cevap}
\newcounter{sor}[section]
\newcommand{\SC}[2]{\stepcounter{sor}\par\noindent\textbf{S\thesection.\thesor\ — #1}\par\nopagebreak
  \noindent\hangindent=1.2em\hangafter=0 #2\par\vspace{0.35em}}

\newcommand{\OC}{O$-$C}
\newcommand{\Pdot}{\ensuremath{\dot{P}}}
\newcommand{\msyil}{\,ms/yıl}
\newcommand{\sR}{\ensuremath{\sigma_{\text{Rømer}}}}

\begin{document}

% ===================================================== KAPAK
\begin{titlepage}
\centering
\vspace*{2cm}
{\LARGE\bfseries TÜBİTAK 2209-A Başvuru Kılavuzu\par}
\vspace{0.5cm}
{\Large Sıcak Jüpiterlerde Yörünge Dönemi Değişimlerinin \OC{} Analiziyle İncelenmesi ve Fiziksel Kökenlerinin Belirlenmesi\par}
\vspace{1.5cm}
{\large Proje yürütücüsü: İdris Avlar\\ Proje ortağı: Bahar Cayib\\ Akademik danışman: Hüseyin Er\\[0.3cm]
Atatürk Üniversitesi, Fen Fakültesi, Astronomi ve Uzay Bilimleri Bölümü\par}
\vspace{1.5cm}
\begin{minipage}{0.85\textwidth}
\begin{bilgi}[Bu kılavuzda ne var?]
\begin{enumerate}
\item \textbf{Adayları nasıl seçtim}: hoca ya da hakem sorarsa söyleyeceklerin, sayılarıyla.
\item \textbf{Başvuru adımları}: TYBS'de adım adım ne yapacağın ve son tarihler.
\item \textbf{Alan alan olası sorular ve cevapları}: formun her bölümünden.
\item \textbf{Şu ana kadar yapılan her şey}: neyi, neden, neye göre yaptığımız.
\item \textbf{Formda olmayan denemeler}: yapıldı ama forma konmadı; sorulursa bil.
\item \textbf{Kavramlar sözlüğü}, \textbf{son kontrol listesi}.
\item \textbf{Ek A}: TYBS'ye yapıştırılacak form metninin tamamı.
\end{enumerate}
\end{bilgi}
\end{minipage}
\vfill
{\small Son güncelleme: 7 Kasım 2026 \quad·\quad Son başvuru: 11 Kasım 2026, saat 17.30 (hedef: 10 Kasım)}
\end{titlepage}

\tableofcontents
\newpage

% =====================================================================
\section{Adayları nasıl seçtim? (Hoca sorarsa)}
% =====================================================================

\begin{soyle}[Tek cümlelik cevap]
``Wang vd. (2024)'nin LOOCV testini geçen 14 adayından, önceden belirlediğim iki veri yeterliliği ölçütüyle seçtim: TESS'te yeterince geçiş zamanı olması (en az 30) ve Rømer etkisinin radyal hız verisiyle sınandığı yayımlanmış bir analizin bulunması. Dört aday ikisini de sağladı: WASP-12 b, WASP-4 b, HAT-P-7 b ve TrES-1 b.''
\end{soyle}

\subsection{Seçim üç adımdan oluşuyor}

\begin{center}
\begin{tabular}{@{}lll@{}}
\toprule
Adım & Ölçüt & Kalan aday \\ \midrule
Başlangıç & Wang vd. (2024)'nin dönem değişimi bildirdiği adaylar & 26 \\
(i) LOOCV & Wang vd.'nin LOOCV testini geçmiş olmak (Tablo 3, ``pass'') & 14 \\
(ii) TESS verisi & Wang vd.'nin TESS geçiş zamanı sayısı $\geq 30$ & 7 \\
(iii) Radyal hız & Rømer'i RV ile sınayan yayımlanmış analiz var & \textbf{4} \\ \bottomrule
\end{tabular}
\end{center}

\subsection{Neden Wang vd. (2024) ve neden LOOCV?}
\begin{itemize}
\item Bu hocanın direktifiydi: ``Wang vd.'nin LOOCV'yi geçen 14 adayından birkaçını gerekçesiyle seç.''
\item Wang vd. (2024), TESS verisiyle 326 sıcak Jüpiteri incelemiş, 26'sında uzun dönemli dönem değişimi bulmuş (18'i kısalma, 8'i uzama).
\item \textbf{LOOCV} (``birini dışarıda bırak çapraz doğrulama''): verilerden her seferinde \emph{bir} geçiş zamanını çıkarıp analizi tekrarlamak. Tek bir ölçüm çıkarılınca sonuç kayboluyorsa (ör. anlamlılık 3σ'nın altına düşüyorsa), o sonuç o tek ölçüme bağlıdır, güvenilmez. 26 adaydan 12'si bu testi geçemedi; 14'ü geçti.
\item \textbf{LOOCV'yi ben yapmadım}; Wang vd.'nin sonucunu (Tablo 3'ün son sütunu) kullandım. Ne olduğunu öğrendim ama proje kapsamında değil.
\end{itemize}

\subsection{Ölçüt (ii): TESS geçiş sayısı en az 30 — neden?}
\begin{itemize}
\item Projenin işi, TESS'in yeni verileriyle dönem değişimini \emph{güncellemek}. Bunun için yıldızın TESS tarafından tekrar tekrar gözlenmiş olması gerekir.
\item Bir TESS sektörü yaklaşık 27 gündür. 14 aday arasında en kısa dönemli gezegende (KELT-16 b, 0,97 gün) bile bir sektörde en fazla $\sim$27 geçiş olur. Yani \textbf{30 geçiş = en az iki sektör}. İki sektör gözlenmiş bir yıldızın sonraki yıllarda yeni sektörlerinin olması da çok olasıdır.
\item \textbf{Nasıl saydım?} Wang vd.'nin VizieR'de yayımladığı geçiş zamanı tablosunda (J/ApJS/270/14) her adayın TESS satırlarını saydım. Sayılar makalenin metniyle aynı (tek fark: WASP-12 b için tabloda 81, metinde 82; sonuç değişmiyor).
\item \textbf{``Neden 30, neden 20 değil?''} sorusunun cevabı: geçenlerin en azı 32, kalanların en çoğu 16 geçişe sahip. Eşik 17 ile 32 arasında \emph{hangi sayı} seçilirse seçilsin aynı 7 aday kalıyor. Yani sonuç eşiğin tam değerine bağlı değil.
\end{itemize}

\subsection{Ölçüt (iii): Radyal hız — neden zorunluydu?}

\begin{dikkat}[En önemli kısım: bunu çok iyi anla]
\OC{} diyagramındaki \textbf{parabol, dönemin değiştiğini gösterir ama nedenini göstermez.} İki farklı neden \OC{}'de birebir aynı parabolü çizer:
\begin{enumerate}
\item \textbf{Gelgit bozulması (gerçek değişim):} yıldızdaki gelgit kabarıklığı gezegeni yavaşlatır, gezegen içeri sarmallanır, yörünge gerçekten küçülür, dönem gerçekten kısalır.
\item \textbf{Rømer etkisi (görünür değişim):} sistemde uzak, görünmeyen bir eşlikçi (yıldız ya da dev gezegen) varsa bütün sistem bakış doğrultusunda ivmelenir. Sistem bize yaklaştıkça her geçişin ışığı bize biraz daha kısa yoldan gelir; geçişler biraz daha erken görünür, dönem kısalıyormuş \emph{gibi} olur. Yörünge ise hiç değişmemiştir.
\end{enumerate}
Geçiş zamanlarına bakarak bu ikisini ayırt edemezsin. \textbf{Ayırt etmenin yolu radyal hızdır (RV):} sistem ivmeleniyorsa yıldızın radyal hızında yıllar boyunca yavaş, sürekli bir kayma (doğrusal eğilim) görülür. Bu eğimden Rømer etkisinin vereceği dönem değişimi hesaplanır:
\[ \dot{P}_\text{Rømer} = \frac{P \cdot (dv/dt)}{c} \]
Bu, \OC{}'den bulunan değerle aynıysa değişim Rømer'dendir (gerçek değil); çok farklıysa Rømer tek başına açıklayamaz, gelgit gibi gerçek bir neden aranır.
\end{dikkat}

\textbf{Literatürden gerçek örnekler (hocaya söyleyebilirsin):}
\begin{itemize}
\item \textbf{WASP-4 b}: aynı dönem kısalmasını Baştürk vd. (2025) gelgitle açıkladı; Bouma vd. (2020) ve Winn ve Stefánsson (2025) RV ile $\sim$8 AB uzaklıkta $\sim$7 Jüpiter kütleli bir eşlikçinin Rømer etkisi olduğunu gösterdi.
\item \textbf{WASP-12 b}: tek kesin örnek sayılmasının nedeni, RV'nin Rømer'i dışlamış olması (Yee vd., 2020).
\end{itemize}

Hocanın direktifinin ikinci yarısı ``değişim varsa nedenini belirle'' idi. Ayırt edilmesi en zor ikili gelgit ile Rømer, ve yalnızca RV ile ayrılabiliyor. Rømer'i RV ile sınanmamış bir aday seçseydik, değişim bulsak bile ``nedeni şu'' diyemezdik.

\begin{soyle}[Hocaya tek cümle]
``\OC{}'deki parabol hem gerçek gelgit bozulmasından hem de uzak bir eşlikçinin ışık-yolu (Rømer) etkisinden kaynaklanabiliyor; ikisi zamanlamada aynı izi bıraktığı için ayırt etmenin tek yolu yıldızın radyal hızındaki uzun dönemli eğilim. Bu yüzden Rømer etkisinin RV ile zaten sınandığı, yani bu ayrımın yapılabileceği adayları seçtim.''
\end{soyle}

\subsection{Rømer analizlerini nereden bulduk? (Hesap yok, hepsi literatür)}
Kendimiz RV ölçmedik ve RV'ye model de uydurmadık. Her aday için, Rømer katkısını RV ile zaten sınamış \textbf{yayımlanmış} bir çalışma aradık:

\begin{center}\small
\begin{tabularx}{\linewidth}{@{}lX@{}}
\toprule
Kaynak & Ne veriyor \\ \midrule
Bouma vd. 2020, Tablo 5 & 51 sıcak Jüpiter için (çoğu Knutson vd. 2014'ün yıllara yayılan Keck/HIRES RV'lerinden) ölçülen RV ivmesi ve buna karşılık gelen Rømer öngörüsü $\dot P_\text{Rømer}\pm\sigma$. WASP-12 b, WASP-4 b, HAT-P-7 b burada. VizieR'de de var: J/ApJ/893/L29/table5. \\
Hagey vd. 2025 & TrES-1 b'nin 12 yıllık RV'sini inceliyor: en iyi model \emph{eğilimsiz}; iddia edilen dönem değişimini açıklayacak bir ivme ``veride fark edilirdi, görülmüyor''. Rømer'i açıkça dışlıyor. \\
Rothmeier vd. 2025 & TrES-5 b için: ``yalnızca 217 güne yayılan 8 RV ölçümü yayımlanmış; uzun dönemli eğilimi belirlemek için kullanışlı değil.'' \\
Oberst vd. 2017 & KELT-16 b keşif makalesi: RV'ler yaklaşık 6 aya yayılıyor; uzun dönemli eğilim için yetersiz. \\ \bottomrule
\end{tabularx}
\end{center}
WASP-46 b için incelediğim çalışmalarda (Wang vd. 2024, Bouma vd. 2020 ve diğerleri) Rømer'i RV ile sınayan bir analiz bulamadım. Formda da ``incelenen literatürde bulunamadı'' yazıyor (``dünyada hiç yok'' demiyoruz).

\subsection{Bouma vd. (2020) Tablo 5 nedir?}
Bouma ve arkadaşları, WASP-4 b'nin dönem kısalmasının Rømer'den olup olmadığını araştırırken aynı hesabı 51 sıcak Jüpiter için yapıp tablo olarak yayımladılar. Her gezegen için:
\begin{itemize}
\item yıldızın RV'sindeki uzun dönemli eğim (ivme, $dv/dt$, m/s/gün),
\item bu ivmenin yaratacağı görünür dönem değişimi $\dot P_\text{Rømer} = P\cdot(dv/dt)/c$ ve belirsizliği (\msyil).
\end{itemize}
Yani tablodaki sayı, ``bu sistemde Rømer etkisi ne kadar dönem değişimi yaratır'' sorusunun yayımlanmış cevabıdır. Biz sadece okuduk.

\subsection{Ölçüt (iii) tam olarak ne?}
Adayın Rømer katkısını RV ile sınayan yayımlanmış bir analiz bulunmalı ve bu analiz ya
\begin{itemize}
\item Rømer öngörüsünü, Wang vd.'nin $|\Pdot|$ değerinin \textbf{üçte birinden küçük} bir belirsizlikle vermeli (WASP-12, WASP-4, HAT-P-7), ya da
\item Rømer etkisini açıkça dışlamalı (TrES-1, Hagey vd. 2025).
\end{itemize}
\textbf{Neden üçte bir?} Rømer'in katkısını, aradığımız sinyalden en az üç kat daha ince bilmek için. Böylece ``Rømer bu değişimi açıklıyor/açıklamıyor'' sonucunu en az 3σ düzeyinde söyleyebiliriz (3σ, ``anlamlı'' için kullanılan standart sınır; Wang vd. de 3σ kullandı).

\textbf{Sonuç eşiğe duyarlı değil}: seçilenlerde belirsizlik, sinyalin yalnızca \%4--6'sı (WASP-12: 1,1/30,2; WASP-4: 0,4/6,4; HAT-P-7: 0,9/18,3). Elenenlerde ise ya hiç analiz yok ya da veri birkaç ayla sınırlı. Eşik 1/3 yerine 1/10 bile olsa aynı dört aday kalır.

\subsection{Aday aday sonuç}
{\small
\begin{longtable}{@{}lrrrp{6.2cm}l@{}}
\toprule
Aday & $|\Pdot|/\sigma$ & TESS & Eşik & Rømer analizi (\msyil) & Sonuç \\ \midrule\endhead
CoRoT-2 b & 7,3 & 5 & 7,2 & -- & (ii) ile elendi \\
HAT-P-37 b & 5,1 & 8 & 11,4 & -- & (ii) ile elendi \\
WASP-10 b & 7,9 & 8 & 9,2 & $-15{,}6\pm4{,}1$ (Bouma 2020) & (ii) ile elendi \\
WASP-22 b & 4,2 & 5 & 23,9 & $+21{,}7\pm2{,}8$ (Bouma 2020) & (ii) ile elendi \\
WASP-45 b & 8,0 & 14 & 56,4 & -- & (ii) ile elendi \\
XO-4 b & 3,6 & 16 & 20,9 & $+4{,}3\pm4{,}3$ (Bouma 2020) & (ii) ile elendi \\
WASP-1 b & 4,3 & 9 & 7,5 & $+0{,}8\pm1{,}5$ (Bouma 2020) & (ii) ile elendi \\
KELT-16 b & 5,8 & 65 & 9,3 & Sınırlayıcı değil: RV'ler $\sim$6 ay (Oberst 2017) & (iii) ile elendi \\
TrES-5 b & 4,6 & 66 & 3,2 & Sınırlayıcı değil: 217 günde 8 RV (Rothmeier 2025) & (iii) ile elendi \\
WASP-46 b & 18,1 & 32 & 17,2 & İncelenen literatürde bulunamadı & (iii) ile elendi \\
\textbf{WASP-12 b} & 32,8 & 81 & 10,1 & $-0{,}1\pm1{,}1$ (Bouma 2020) & \textbf{Seçildi} \\
\textbf{WASP-4 b} & 11,5 & 47 & 2,1 & $-5{,}9\pm0{,}4$ (Bouma 2020) & \textbf{Seçildi} \\
\textbf{HAT-P-7 b} & 4,5 & 66 & 6,1 & $+15{,}0\pm0{,}9$ (Bouma 2020) & \textbf{Seçildi} \\
\textbf{TrES-1 b} & 6,3 & 34 & 4,9 & Dışlandı: 12 yıllık RV'de eğilim yok (Hagey 2025) & \textbf{Seçildi} \\ \bottomrule
\end{longtable}}
\vspace{-0.5em}
{\footnotesize $|\Pdot|/\sigma$: Wang vd.'nin dönem değişim hızının belirsizliğine oranı (Tablo 3, 5σ kırpma sütunu). Eşik = $|\Pdot_\text{Wang}|/3$, \msyil. Bouma vd.'nin asimetrik hataları (ör. $+1{,}11/-1{,}07$) yuvarlanarak tek değerle yazıldı. Bouma'da değeri olan dört aday (WASP-10, -22, -1, XO-4) zaten (ii)'de elendiği için o değerler seçimi etkilemiyor; tablo dolu olsun diye yazıldı.}

\textbf{Bu sayılar ne söylüyor? (Sorulursa)} WASP-12'de Rømer öngörüsü sıfıra çok yakın ($-0{,}1\pm1{,}1$), oysa gözlenen değişim $\sim-30$: Rømer açıklayamaz, gelgit güçlü aday. WASP-4'te Rømer öngörüsü ($-5{,}9$) gözlenen değişime ($\sim-6$) çok yakın: Rømer büyük ölçüde açıklayabilir. HAT-P-7'de Rømer $+15$, gözlenen $\sim+18$: önemli kısmı Rømer olabilir. Bunlar zaten H1 ve H2 hipotezlerinin dayanağı. Ama kesin hükmü projede (İP4) vereceğiz.

\subsection{Seçilen dört hedef neden iyi bir set?}
Rastgele değil, birbirini tamamlayan dört durumu temsil ediyorlar:
\begin{itemize}
\item \textbf{WASP-12 b}: gelgit bozulması kesinleşmiş tek örnek $\rightarrow$ yöntemin doğru çalıştığını göstermek için \emph{kalibrasyon}.
\item \textbf{WASP-4 b}: gelgit mi Rømer mi tartışmalı $\rightarrow$ asıl \emph{ayırt etme} testi.
\item \textbf{HAT-P-7 b}: dönem \emph{uzuyor}; yıldız yörüngeden yavaş döndüğü için gelgit dönemi uzatamaz $\rightarrow$ gelgit dışı bir neden (Rømer / eşlikçi yıldız, Yang vd. 2025).
\item \textbf{TrES-1 b}: farklı çalışmalar farklı değerler buluyor ($-16$, $-14{,}9$, $-7{,}1$\msyil) $\rightarrow$ değişim gerçekten var mı, yoksa literatür zamanlarının kalitesinden mi?
\end{itemize}

\subsection{Bu seçimde ne yaptım, ne yapmadım? (Dürüst cevap)}
\begin{soyle}[Doğru ifade]
``Ön çalışmada hesap ya da test yapmadım; literatür taraması yaptım. 14 adayın TESS geçiş sayılarını Wang vd.'den, Rømer etkisinin RV ile sınanıp sınanmadığını yayımlanmış analizlerden (Bouma vd. 2020, Hagey vd. 2025, Rothmeier vd. 2025, Oberst vd. 2017) inceledim ve iki ölçütü sağlayan dört hedefi seçtim.''
\end{soyle}
\begin{itemize}
\item Tablo 1'deki \textbf{her sayı} bir makaleden okunuyor ya da basit bir bölme (Eşik = $|\Pdot|/3$). Kendi hesabın yok.
\item \textbf{Adayları kendi \OC{} sonuçlarıma bakarak seçmedim.} Ölçütler projenin sonucundan bağımsız; ``beğendiğim sonuçları seçmişim'' denemez.
\item Tablo 1'in ham hali ve kaynakları: \texttt{veri/tablo1\_secim.csv}.
\end{itemize}

\subsection{Olası itirazlar}
\SC{Neden sadece dört hedef?}{Sayıyı ben belirlemedim; ölçütler 14 adayı dörde indirdi. Ayrıca 12 aylık bir lisans projesi için dört hedef, her birinin altı fiziksel etkiye göre ayrı ayrı incelenmesine izin veren gerçekçi bir sayı.}
\SC{Rømer'i sınanmamış adaylarda da dönem değişimini ölçebilirdin, neden seçmedin?}{Ölçebilirdim ama nedenini belirleyemezdim; proje ``değişim varsa nedenini belirle'' diyor. Rømer sınanmamış bir adayda gelgit ile Rømer ayırt edilemez, ve tayfçekerimiz olmadığı için yeni RV de alamayız.}
\SC{Neden LOOCV'yi geçemeyen 12 adayı almadın?}{Wang vd.'nin kendisi onları ``zayıf aday'' sayıyor: sonuçları tek bir ölçüme bağlı. Hocanın direktifi de geçen 14'ten seçmekti.}
\SC{TrES-1 için sayı yok, neden seçildi?}{Hagey vd. (2025) 12 yıllık RV'yi incelemiş ve eğilim bulmamış; dönem değişimini açıklayacak bir ivmenin veride görüneceğini, görünmediğini açıkça yazıyor. Yani Rømer'i dışlıyor; ölçütün ikinci yolu bu.}
\SC{KELT-16 b'nin RV verisi yok muydu?}{Var, keşif makalesinde (Oberst vd., 2017) 20 ölçüm var ama hepsi yaklaşık 6 ay içinde alınmış. Yıllara yayılan bir eğilimi ölçmeye yetmez; Rømer sınanmış sayılmaz.}
\SC{WASP-46 b için gerçekten hiç analiz yok mu?}{Benim incelediğim çalışmalarda yok. ``Dünyada hiç yok'' demiyorum, formda da ``incelenen literatürde bulunamadı'' yazıyor.}
\SC{Seçimi senin O$-$C sonuçların etkiledi mi?}{Hayır. Ölçütler yalnızca veri yeterliliğine bakıyor.}

% =====================================================================
\section{Başvuru adımları}
% =====================================================================

\begin{dikkat}[Tarihler]
Son başvuru: \textbf{11 Kasım 2026, saat 17.30} (sistem mesai bitiminde kapanır). Sen \textbf{10 Kasım'da} bitir ki hocanın onaylamaya ve üniversitenin e-imzasına vakti kalsın. Başvuru, kuruluş yetkilisi e-imzası atılınca geçerli sayılır.
\end{dikkat}

\begin{bilgi}[Not]
TYBS ekranlarını ben görmedim; aşağıdaki adımlar çağrı metninden ve form bilgi notundan çıkarıldı. Ekrandaki başlıkların adı biraz farklı olabilir; sırası ve içeriği aynıdır.
\end{bilgi}

\subsection{Başvurudan önce (bugün/yarın)}
\begin{enumerate}
\item \textbf{ARBİS} (arbis.tubitak.gov.tr): kendi kaydın eksiksiz ve güncel olsun. Bahar'ınki de.
\item \textbf{Hocanın bilgileri} YÖKSİS ve ARBİS'te güncel olmalı (sistem oradan kontrol ediyor). Hocana bunu hatırlat.
\item \textbf{Bahar'ın koşulları}: Türkiye'de lisans/ön lisans öğrencisi (açıköğretim ve hazırlık olamaz), ARBİS kaydı var, bu dönem başka 2209-A/B başvurusunda ya da süren projede yok. Sen de aynı koşulları sağlıyorsun.
\item Hocana Word belgesini gönder, onayını al.
\end{enumerate}

\subsection{TYBS'de (tybs.tubitak.gov.tr)}
\begin{enumerate}
\item Giriş yap, 2209-A 2026 çağrısına yeni başvuru oluştur.
\item \textbf{Genel bilgiler}: proje başlığı (Ek A'dan, aynen), akademik danışman (Hüseyin Er), proje ortağı (Bahar Cayib). Bilimsel alan seçimi istenirse astronomi/astrofizik ile ilgili olanı seç.
\item \textbf{Bütçe}: \textbf{0 TL}. Gerekçe istenirse: ``Proje yalnızca açık arşiv verisi ve ücretsiz yazılımla yürütülecek; ekipman, sarf, seyahat veya hizmet alımı gerekmemektedir.''
\item \textbf{Araştırma önerisi metin kutuları}: Ek A'daki metni bölüm bölüm ilgili kutuya yapıştır: Özet, Anahtar Kelimeler, 1.1, 1.2, 2, 3.1 (takvim), 3.2 (risk), 3.3 (olanaklar), 4 (yaygın etki), 5 (diğer konular + Tablo 1).
\item \textbf{Tabloları} metin kutusunun üstündeki \textbf{``Tablo''} düğmesiyle kur (yapıştırınca bozulabilir). Takvim, risk, olanaklar, yaygın etki ve Tablo 1 için geçerli.
\item \textbf{Özel karakterler} (σ, Δ, ±, ≥, ′, ⁵): Yunan harfleri doğrudan yapıştırılabilir; görünmezse ``Özel karakter'' düğmesini kullan.
\item \textbf{Kaynakça adımı}: 35 kaynağı ayrı ayrı gir (Ek A'nın sonu). Hepsinin sonunda DOI var (DOI zorunlu); DACE, MAST ve iki resmi belge için bağlantı ve erişim tarihi var.
\item \textbf{Yapay zekâ beyanı}: kendin işaretle ve yaz.
\item \textbf{Proje Önerisini Önizle}: her bölümün göründüğünü, tabloların taşmadığını kontrol et. Boş görünen kısım olursa bilgi notunun dediği gibi biçimi temizleyip yeniden yapıştır.
\item \textbf{Onayla} $\rightarrow$ başvuru salt okunur olarak hocanın e-postasına ``danışman onayı'' için gider. Hoca onaylarsa üniversitenin kuruluş yetkilisine e-imza için gider. Hoca iade ederse düzeltip yeniden gönderirsin.
\end{enumerate}

\subsection{Ön incelemede elenmemek için}
\begin{itemize}
\item Özet 25–450 kelime (bizimki 275). 1.1, 1.2 ve 2 en az 350 karakter (hepsi çok üstünde).
\item Takvimdeki başarı yüzdeleri toplamı 100 (20+20+25+35). Toplam süre $\leq$ 12 ay.
\item Literatür taraması, rapor yazma, sonuç paylaşımı ayrı iş paketi olamaz (bizde yok).
\item Hiçbir kurum, firma ya da marka öne çıkarılmamalı (madde 4.1.8). Bu yüzden forma dergi adı örneği ya da kod deposu sitesi adı yazmadık.
\item Her kaynağa metinde atıf yapılmalı (yapıldı, kontrol edildi).
\end{itemize}

\subsection{Başvurudan sonra}
\begin{itemize}
\item Sonuçlar e-bideb.tubitak.gov.tr'den duyurulur. Değerlendirme ağırlıkları: bilimsel nitelik \%35, yöntem \%25, proje yönetimi \%20, yaygın etki \%20.
\item Etik kurul gerekmez (insan/hayvan deneği yok, yalnızca açık astronomi verisi).
\item Kabul edilirse proje en fazla 12 ay sürer; sonuç raporu TÜBİTAK'ın formatında yüklenir. Üretilen her çıktıda (ör. makale) TÜBİTAK desteği belirtilir.
\end{itemize}

% =====================================================================
\section{Alan alan olası sorular ve cevaplar}
% =====================================================================
Her cevabı ezberleme; \emph{mantığını} anla. Kısa cevap ver, sorulursa ayrıntıya gir.

\subsection{Başlık ve konu}
\SC{Projenin konusu ne, bir cümleyle?}{Dört sıcak Jüpiterin yörünge döneminin değişip değişmediğini yeni TESS verisiyle \OC{} diyagramlarından sınamak ve değişim varsa bunun gelgit, apsis hareketi, Rømer, üçüncü cisim, Applegate ya da kütle kaybından hangisinden kaynaklandığını belirlemek.}
\SC{Sıcak Jüpiter nedir?}{Jüpiter kütlesinde, yıldızının etrafında birkaç günde dolanan gezegen. Yıldıza çok yakın oldukları için gelgit etkileri güçlüdür.}
\SC{Bu neden önemli?}{Dönem değişimi, yıldızın gelgit enerjisini ne kadar verimli soğurduğunu ($Q'_*$) doğrudan ölçmenin ender yollarından biri. Kuramsal $Q'_*$ tahminleri birkaç mertebe farklı (Penev vd., 2018); gözlem bu belirsizliği daraltır. Ayrıca gezegenlerin sonunda yıldızlarına düşüp düşmeyeceğini söyler.}

\subsection{Özet}
\SC{Özette ``bütçe gerektirmeyen'' diyorsun, neden?}{Tüm veriler açık arşivlerde, yazılımlar ücretsiz, gözlem yok. Harcanacak kalem yok; uydurma kalem yazmak tutarlılığı bozar.}
\SC{Çıktılar ne?}{Dört hedef için güncel efemeris (T0, P, varsa dP/dt), her hedef için fiziksel köken hükmü, açık erişimli veri seti ve analiz kodu; yeterli görülürse danışmanla kısa bir makale.}

\subsection{1.1 Konunun önemi ve bilimsel nitelik}
\SC{Gelgit bozulması nasıl oluyor?}{Gezegen yıldızda bir gelgit kabarıklığı oluşturur. Yıldız, gezegenin yörüngede dolanmasından yavaş dönüyorsa kabarıklık gezegenin gerisinde kalır ve gezegeni geriye çeker. Gezegen açısal momentum kaybeder, içeri sarmallanır, dönem kısalır.}
\SC{$Q'_*$ nedir?}{Değiştirilmiş yıldız gelgit kalite faktörü. Küçükse yıldız gelgit enerjisini verimli soğurur, bozulma hızlıdır. WASP-12 için $\approx 1{,}8\times10^5$ (Yee vd., 2020).}
\SC{Literatürdeki eksik ne, sen neyi gideriyorsun?}{Wang vd. adayları buldu ama fiziksel kökenlerini sistem sistem sınamadı. Ayrıca \OC{}'deki parabol tek başına kökeni göstermiyor; WASP-4 ve TrES-1 gibi örneklerde farklı gruplar farklı sonuçlar buluyor. Ben (i) Wang sonrası TESS sektörleriyle zaman tabanını $\sim$2--6 yıl uzatacağım, (ii) altı etkiyi her hedefte aynı karar ağacıyla değerlendireceğim, (iii) gelgit ile Rømer'i arşiv RV'siyle ayıracağım.}
\SC{Bu gezegenler zaten çalışılmamış mı?}{Çalışılmış, ve tam da bu yüzden seçildiler: sonuçlar çelişkili (WASP-4: gelgit mi Rømer mi; TrES-1: $-16$, $-14{,}9$, $-7{,}1$\msyil). WASP-12 ise yöntemin kalibrasyonu. Yeni TESS verisi ve aynı yöntemle sistematik inceleme özgün kısım.}
\SC{``Zaman tabanı 2--6 yıl uzayacak'' nereden?}{Wang vd.'nin son TESS zamanından bugüne kadar çıkan sektörlerin tarihlerinden: WASP-12, HAT-P-7 ve TrES-1 için $\sim$2 yıl, WASP-4 için $\sim$6 yıl.}
\SC{Araştırma sorusu ne?}{Seçilen dört sıcak Jüpiterin dönemi güncel TESS verisiyle anlamlı biçimde değişiyor mu; değişiyorsa sorumlusu gelgit, apsis, Rømer, üçüncü cisim, Applegate, kütle kaybından hangisi?}
\SC{Hipotezlerin ne, neden böyle?}{H1: WASP-12 b'de değişim literatürle uyumlu, Rømer açıklayamaz (gelgit) — çünkü Yee vd. ve Winn ve Stefánsson bunu gösterdi. H2: WASP-4 b ve HAT-P-7 b'deki değişimin önemli kısmı Rømer — çünkü ikisinin RV'sinde belirgin eğilim var (Bouma vd. 2020; Winn vd. 2009). H3: TrES-1 b'deki kısalma gerçekçi hatalarla anlamlı değil — çünkü değerler çalışmadan çalışmaya çok değişiyor ve Adams vd. literatür zamanlarında hatalar buldu.}
\SC{Hipotezin tutmazsa ne olur?}{O da bir sonuç; raporlanır. Proje ``değişim var mı?'' sorusunu soruyor, ``var'' çıkmak zorunda değil.}
\SC{Kalkınma Planı ile ilişki ne?}{On İkinci Kalkınma Planı madde 270 temel bilimlere özel önem verilmesini ve uzay araştırmalarını kritik alan saymayı, madde 271 genç nesillerin uzay alanında çalışmaya teşvik edilmesini ve uzay araştırmalarının yaygınlaştırılmasını söylüyor. 2030 Sanayi ve Teknoloji Stratejisi de TÜBİTAK araştırma desteklerinin ``Fizik ve Astronomi'' dahil temel bilimlerde artırılmasını öngörüyor. Proje bir uzay görevinin (TESS) verisiyle lisans düzeyinde araştırmacı yetiştiriyor. (Çağrı metnine göre bu alanlarla ilişkili projeler için destek taban puanı ayrı belirlenebiliyor.)}

\subsection{1.2 Amaç ve hedefler}
\SC{Hedef 1'deki ``her hedef için en az 70 zaman'' nereden?}{Katalog ve literatür zamanları ile yeni TESS zamanları birlikte. Denemede dört hedef için 94 ile 227 arasında zaman toplandı; 70 güvenli bir alt sınır.}
\SC{``En az 150 yeni zaman, ortanca belirsizlik $\leq$ 60 s'' gerçekçi mi?}{Evet. Denemede Wang sonrası sektörlerden 174 zaman ölçüldü, hedef başına ortanca hata 31--56 s çıktı.}
\SC{Hedef 2'de ``WASP-12 b'de literatürle 2σ uyum'' neden?}{WASP-12 b'nin değeri iyi biliniyor ($-29\pm2$\msyil). Benim yöntemim de buna yakın bir şey bulursa yöntemin doğru çalıştığını gösterir; bulmazsa bir hata vardır.}
\SC{Hedef 3'teki hükümler ne demek?}{Her etki için: dışlandı (veri açıkça reddediyor), olası değil, olası, en olası, ayrılamadı (veri ayırt etmeye yetmiyor). Önceden yazılmış kurallarla verilecek ki keyfî olmasın.}
\SC{Efemeris nedir, neden yayımlanacak?}{Bir gezegenin geçiş zamanlarını öngören formül: $T = T_0 + P\cdot E$ (gerekirse $+\tfrac12\,\tfrac{dP}{dE}E^2$). Gözlemciler bir sonraki geçişi ne zaman gözleyeceklerini bununla planlar.}

\subsection{2 Yöntem — genel}
\SC{Araştırma tasarımın ne?}{Açık arşiv verisine dayanan, hipotez sınayan gözlemsel bir çalışma. Bağımsız değişken geçiş epoku ($E$), bağımlı değişken orta geçiş zamanı ve ondan türetilen \OC; türetilen asıl büyüklük dP/dt. Kontrol edilen değişkenler: zaman sistemi, ölçüm belirsizliği, veri kaynağı.}
\SC{Hangi arşivleri kullanacaksın?}{TESS ışık eğrileri: MAST. Katalog tabloları: VizieR. Sistem parametreleri: NASA Exoplanet Archive. RV: DACE ve makale tabloları. Yer gözlemi zamanları: ETD/VarAstro ve ExoClock. Dört hedefin hepsinin bu arşivlerde kaydı olduğu kontrol edildi (ETD'de de, ExoClock'ta da).}

\subsection{İP1 — Zaman tablosu (Ay 1--2)}
\SC{BJD\_TDB nedir, neden dönüştürüyorsun?}{Zamanı Güneş Sistemi'nin kütle merkezine taşınmış ve düzgün akan bir saate (TDB) göre ifade eden sistem. Eski makaleler HJD ya da BJD\_UTC kullanabiliyor. Farklar dakikalara varabiliyor ve sahte dönem değişimi yaratabiliyor (Adams vd. 2024 bunun örneklerini buldu). Hepsini tek sisteme (BJD\_TDB) çeviriyoruz.}
\SC{Zaman sistemi belirlenemezse?}{Bu zamanlar işaretlenecek, analiz onlar dahil ve hariç iki kez yapılacak; sonuç değişiyorsa raporlanacak.}
\SC{ETD'de ``yüksek kalite notlu'' ne demek?}{ETD her ışık eğrisine 1--5 arası kalite notu (DQ) veriyor. Yalnızca iyi notlu, tam geçişleri kullanacağız.}
\SC{``Epok numarasının kesirli kısmı 0,02'yi aşan zamanlar'' ne?}{Bir zamanın epoku tam sayıya çok yakın olmalı. Kesirli kısım büyükse ya zaman yanlış yazılmış ya da zaman sistemi farklı demektir; bunlar ayrıca kontrol edilecek.}

\subsection{İP2 — Yeni TESS zamanları (Ay 3--4)}
\SC{Orta geçiş zamanını nasıl ölçeceksin?}{TESS'in 2 dakikalık ışık eğrileri indirilip normalize edilecek. Önce tüm geçişler katlanıp geçişin şekli (gezegen/yıldız yarıçap oranı, a/R*, eğiklik) batman modeliyle bulunacak. Sonra her geçişe ayrı ayrı, yalnızca orta zaman ve taban çizgisi serbest bırakılarak model uydurulacak. Orta zamanın belirsizliği emcee ile MCMC örneklemesinden alınacak.}
\SC{MCMC nedir, neden?}{Parametrelerin olası değerlerini rastgele ama akıllıca gezerek dağılımını çıkaran yöntem. Belirsizliği gerçekçi verir.}
\SC{Kenar kararması neden sabit?}{Geçiş simetrik olduğu için orta zamanı neredeyse etkilemez; serbest bırakmak yalnızca hesabı yavaşlatır ve kararsızlaştırır.}
\SC{Yöntemin doğru olduğunu nasıl bileceksin?}{Wang vd.'nin zaten ölçtüğü bir sektörü yeniden ölçüp ortak geçişlerde karşılaştıracağım. Başarı ölçütü: ortanca fark en fazla 1σ. (Denemede sektör 20 için ortanca fark 0,40σ çıktı; Bölüm 5.)}

\subsection{İP3 — \OC{} analizi (Ay 5--6)}
\SC{\OC{} nedir?}{Gözlenen geçiş zamanı (O) ile sabit dönemli efemerisin hesapladığı zaman (C) arasındaki fark. Dönem sabitse \OC{} düz bir çizgidir; dönem sabit hızla değişiyorsa parabol olur.}
\SC{Neden üç model (doğrusal, kuadratik, kübik)?}{Doğrusal: dönem sabit. Kuadratik: dönem sabit hızla değişiyor (dP/dt). Kübik: değişim hızının kendisi de değişiyor mu (ör. uzun dönemli bir eşlikçinin yörüngesinin bir parçası mı)?}
\SC{dP/dt nasıl bulunuyor?}{Kuadratik modeldeki $\tfrac12\,\tfrac{dP}{dE}E^2$ teriminden $dP/dE$ bulunur, $dP/dt = (1/P)\,dP/dE$; birim \msyil{} yapılır.}
\SC{BIC ve ΔBIC nedir?}{BIC $=\chi^2 + k\ln n$: uyumun iyiliğini ölçer ama fazla parametreyi cezalandırır. İki modelin BIC farkı (ΔBIC) 10'dan büyükse karmaşık model ``çok güçlü'' tercih edilir. Wang vd. de bunu kullandı.}
\SC{Neden hem $|\Pdot|/\sigma\geq3$ hem ΔBIC $>10$?}{Biri değerin sıfırdan ne kadar uzak olduğunu, diğeri parabol modelinin düz çizgiden gerçekten iyi olup olmadığını söyler. İkisi birlikte daha sağlam. Wang vd. ile aynı ölçütler, böylece sonuçlar karşılaştırılabilir.}
\SC{5σ kırpma nedir?}{Modelden 5σ'dan fazla sapan nokta aykırı sayılıp atılır, model yeniden uydurulur; aykırı kalmayana kadar tekrarlanır. 5σ, sadece gerçekten bozuk noktaları atacak kadar geniş.}
\SC{Hata çubuklarını neden ölçekliyorsun?}{Bazı literatür zamanlarının hataları gerçekçi olmayacak kadar küçük ve sahte anlamlılık yaratıyor. Literatür ve TESS gruplarının hataları ayrı ayrı, model artıklarının indirgenmiş $\chi^2$'si 1 olacak şekilde \emph{yalnızca büyütülerek} ölçekleniyor. Hiçbir zaman küçültülmüyor.}

\subsection{İP4 — Fiziksel köken (Ay 7--11)}
\SC{Altı etkiyi tek tek nasıl sınayacaksın?}{(a) \textbf{Rømer}: arşiv RV'sinden eğim $\rightarrow$ $\dot P = P(dv/dt)/c$, \OC{} değeriyle karşılaştır: fark $<2σ$ ise Rømer açıklayabilir, 2--3σ kısmen, $>3σ$ tek başına açıklayamaz. (b) \textbf{Gelgit}: gereken $Q'_*$'yi $\dot P = -\frac{27\pi}{2Q'_*}\frac{M_p}{M_*}\left(\frac{R_*}{a}\right)^5$ ile hesapla, kuramsal aralıkla kıyasla; dönem uzuyorsa gelgit dışlanır. (c) \textbf{Apsis hareketi}: dışmerkezliğin üst sınırından beklenen genlik ($\approx eP/\pi$); WASP-12 b'de örtülme zamanlarının kayma yönü. (d) \textbf{Applegate ve üçüncü cisim}: artıklarda periyodik sinyal (Lomb--Scargle), kübik modelle eğriliğin sabitliği, yıldız etkinliği. (e) \textbf{Kütle kaybı}: gezegenin Roche sınırına uzaklığı ($a/a_\text{Roche}$).}
\SC{Apsis hareketi nedir?}{Hafif dışmerkez bir yörüngenin uzun ekseninin yavaşça dönmesi. Geçiş zamanlarında sinüs benzeri bir dalga yapar; kısa bir zaman aralığında parabol gibi görünebilir. Örtülmeler (gezegen yıldızın arkasına geçerken) geçişlere göre ters yönde kayar; gelgit bozulmasında ise aynı yönde. WASP-12'de bu test yapılabiliyor.}
\SC{Applegate etkisi nedir?}{Yıldızın manyetik etkinlik döngüsü yıldızın şeklini (kütle dağılımını) değiştirir, bu da gezegenin dönemini çevrimsel olarak değiştirebilir. Periyodik bir \OC{} sinyali bekleriz (Watson ve Marsh, 2010).}
\SC{Rømer ile ``üçüncü cismin dinamik etkisi'' farkı ne?}{Rømer: eşlikçi sistemi ivmelendirir, ışığın yolu değişir; yörünge değişmez. Dinamik etki: eşlikçinin çekimi gezegenin yörüngesini gerçekten değiştirir.}
\SC{Rossiter--McLaughlin etkisi nedir, neden atıyorsun?}{Gezegen yıldızın önünden geçerken yıldızın dönen yüzeyinin bir kısmını örter ve ölçülen RV'yi bozar. Bu ölçümler eğilimi yanıltmasın diye çıkarılıyor.}
\SC{RV'ler farklı aletlerden, nasıl birleştiriyorsun?}{Her alete ayrı sıfır noktası (γ) veriyorum; aynı eğimi ortak uyduruyorum. Sonucun hangi veri setine duyarlı olduğunu ayrıca raporlayacağım.}
\SC{Lomb--Scargle nedir?}{Düzensiz aralıklı veride periyodik sinyal arayan yöntem. Yanlış alarm olasılığı 0,01'in altındaki tepeler sinyal sayılır; gözlem mevsiminden gelen $\sim$1 yıllık tepeler sayılmaz.}

\subsection{Ön çalışma}
\SC{Ön çalışmada ne yaptın?}{Literatür taraması: 14 adayın TESS geçiş sayısını Wang vd.'den aldım, her aday için Rømer etkisini RV ile sınayan yayımlanmış analiz aradım (Bouma vd. 2020 Tablo 5, Hagey vd. 2025, Rothmeier vd. 2025, Oberst vd. 2017), iki ölçütü sağlayan dört hedefi seçtim ve bu hedeflerin geçiş zamanlarının ve yeni TESS verilerinin arşivlerde bulunduğunu kontrol ettim. Hesap ya da test yok.}
\SC{Neden grafik koymadın?}{Ön çalışmayı veri yeterliliği ile sınırlı tuttum; ölçüm sonuçlarını projenin kendisinde, yöntemin tamamıyla vermek istedim.}

\subsection{3.1 Çalışma takvimi ve ekip}
\SC{Neden 11 ay iş paketi, proje 12 ay?}{12. ay tampon (gecikmeler için) ve sonuç raporu için bırakıldı; kurallara göre rapor yazımı ayrı iş paketi olamaz.}
\SC{Yüzdeler (20/20/25/35) neye göre?}{İşin ağırlığına göre. En çok emek ve en önemli bilimsel çıktı İP4'te (fiziksel köken), o yüzden \%35.}
\SC{Bahar ne yapacak?}{Her aşamada destek: İP1'de literatür ve arşiv zamanlarının derlenmesi, İP2'de TESS verilerinin indirilmesi ve geçişlerin görsel kontrolü, İP3'te \OC{} diyagramlarının hazırlanması ve veri kontrolü, İP4'te RV ve parametre derlemesi. Modelleme ve fiziksel testler bende.}
\SC{Hoca ne yapacak?}{Akademik danışman: yöntem kararları, sonuçların yorumu, düzenli görüşmelerle ilerleme takibi.}

\subsection{3.2 Risk yönetimi}
\SC{En büyük risk ne?}{Literatür zamanlarındaki zaman sistemi hataları ve gerçekçi olmayan küçük hatalar sahte sinyal üretebilir. Önlem: her kaynağın zaman sistemini kontrol, belirsizlerle/belirsizler olmadan iki analiz, hata ölçekleme.}
\SC{TESS geçişleri ölçülemezse?}{B planı: sektördeki tüm geçişler birlikte uydurulup sektör başına tek zaman alınır; o da olmazsa Wang vd. ve Ivshina ve Winn TESS zamanlarıyla devam edilir. Sadece bu adım değişir, sonraki analiz aynı kalır.}
\SC{Hiç değişim bulamazsan proje boşa mı gider?}{Hayır; bu araştırma sorusunun bir cevabıdır, güncel sabit efemeris verilir.}
\SC{RV ayırt etmeye yetmezse?}{Literatürdeki eğimler ve RV gerektirmeyen testler (Q′*, eğrilik, periyodogram, Roche) kullanılır; gelgit ile Rømer ayrılamazsa ``ayrılamadı'' hükmü gerekçesiyle verilir.}

\subsection{3.3 Araştırma olanakları}
\SC{Bölümde bilgisayar yok, nasıl yapacaksın?}{Kendi dizüstü bilgisayarım ve üniversite kütüphanesindeki bilgisayarlar yeterli. Veri küçük: dört hedefin bütün TESS ışık eğrileri $\sim$76 MB; denemedeki hesaplar sıradan bir bilgisayarda yaklaşık yarım saat sürdü.}

\subsection{4 Yaygın etki}
\SC{Makale yazabilecek misin?}{Hedef, sonuçlar yeterliyse danışmanla birlikte kısa bir makale. Garanti değil; formda da ``yeterli görülmesi halinde'' yazıyor.}
\SC{Veri seti kimin işine yarar?}{Dört hedefin BJD\_TDB'de temizlenmiş tüm geçiş zamanları ve güncel efemerisleri, gelecekte bu gezegenleri gözleyecek herkesin işine yarar.}
\SC{Yeni projeye çıktı?}{Aynı yöntem Wang vd.'nin diğer adaylarına uygulanabilir; bir lisansüstü tez çalışmasına temel olabilir.}

\subsection{5 Diğer konular (Tablo 1)}
\SC{Tablo 1'i nereden aldın?}{Literatürden derledim. Satırlar Wang vd.'nin LOOCV'yi geçen 14 adayı; $|\Pdot|/\sigma$ ve TESS sayısı Wang vd.'nin VizieR tablolarından; Rømer sütunu Bouma vd. 2020 Tablo 5, Hagey vd. 2025, Rothmeier vd. 2025 ve Oberst vd. 2017'den. Bkz. Bölüm 1.}
\SC{``Eşik'' sütunu ne?}{Wang vd.'nin $|\Pdot|$ değerinin üçte biri. Yayımlanmış Rømer öngörüsünün belirsizliği bundan küçükse (ya da Rømer açıkça dışlanmışsa) aday seçiliyor.}

\subsection{6 Kaynakça}
\SC{Kaynaklar doğru mu?}{35 kaynağın hepsinin DOI'si çözülüyor ve yazar, yıl, başlık, dergi, cilt, sayfa bilgileri Crossref kayıtlarıyla eşleşiyor (7 Kasım'da kontrol edildi). Metinde hepsine atıf var.}
\SC{En önemli 5 kaynak ve ne söyledikleri?}{Wang vd. 2024 (aday listesi, LOOCV); Yee vd. 2020 (WASP-12 kesin bozulma, $-29\pm2$\msyil, Rømer dışlandı); Bouma vd. 2020 (WASP-4 RV'si ivmeleniyor, Rømer $\approx -5{,}9$\msyil); Winn ve Stefánsson 2025 (WASP-4 bozulmuyor, $\sim$7 Jüpiter kütleli eşlikçi; tek kesin örnek WASP-12); Patra vd. 2017 (gelgit formülü, WASP-12).}
\SC{Neden bir resmi belgeye (Kalkınma Planı) atıf var?}{Form, projenin bu planlarla ilişkisini açıklamayı istiyor; atıf yapılan her şey kaynakçada olmalı.}

\subsection{Genel / beklenmedik sorular}
\SC{Neden TESS?}{Tüm gökyüzünü tarayan, aynı yıldızı yıllar arayla tekrar gözleyen uzay teleskobu; hassas ve homojen geçiş zamanları veriyor. Wang vd. de TESS kullandı.}
\SC{Kendi gözlemini neden yapmıyorsun?}{Gerek yok: TESS'in yeni sektörleri ve yer gözlemi veritabanları yeterli veri sağlıyor. Ayrıca RV için gereken tayfçeker bölümde yok.}
\SC{LOOCV'yi sen de yapacak mısın?}{Proje kapsamında değil. Wang vd.'nin sonucunu seçim için kullandım.}
\SC{Kod kimin?}{Python'da yazılmış analiz betikleri (astropy, astroquery, batman, emcee, numpy). Projede bunları kendim geliştirip kullanacağım; açık erişimli paylaşılacak.}
\SC{Yapay zekâ kullandın mı?}{Kendi beyanına göre cevapla (formdaki beyanla tutarlı olsun).}

% =====================================================================
\section{Şu ana kadar yapılan her şey}
% =====================================================================

\subsection{Kronolojik özet}
\begin{enumerate}
\item \textbf{Hocanın direktifi alındı}: Wang vd. (2024)'nin LOOCV'yi geçen 14 adayından birkaçını gerekçesiyle seç; LOOCV'yi kendin öğren (proje kapsamı değil); \OC{} diyagramı çiz; TESS ve VarAstro/literatür verisiyle dönemin değişip değişmediğine bak; değişiyorsa gelgit/apsis/Rømer/Applegate/üçüncü cisim/kütle kaybından hangisi olduğunu belirle.
\item \textbf{Kısıtlar belirlendi}: tayfçeker yok $\rightarrow$ RV yalnızca arşivden; bütçe 0 TL; yalnızca açık veri; hocanın söylediklerine fazla ekleme yapma.
\item \textbf{Literatür okundu}: Wang 2024, Ivshina ve Winn 2022, Yee 2020, Patra 2017, Bouma 2020, Baştürk 2025, Winn ve Stefánsson 2025, Adams 2024, Hagey 2025, Kutluay 2026, Yang 2025 vb. Amaç: hangi adayda ne tartışıldığını görmek.
\item \textbf{Analiz kodları yazıldı ve denendi} (Bölüm 5): \OC{} kodu, TESS zaman ölçüm kodu, RV/Rømer kodu. Yöntemin uygulanabilir olduğunu görmek için. Bu kodları yapay zekâ (Claude) yazdı ve çalıştırdı.
\item \textbf{Seçim ölçütleri kesinleşti ve Tablo 1 üretildi} (Bölüm 1). Ölçütler veri yeterliliğine dayandırıldı ki seçim sonuçlardan bağımsız olsun. İlk sürümde Rømer sütunu arşiv RV'lerinden kodla hesaplanmıştı (\sR); senin isteğinle tamamen yayımlanmış analizlere (Bouma vd. 2020 Tablo 5, Hagey vd. 2025, Rothmeier vd. 2025, Oberst vd. 2017) dayandırıldı. Seçilen dört hedef değişmedi.
\item \textbf{Form metinleri yazıldı}, TÜBİTAK kurallarına göre (bilgi notu, çağrı metni) bölüm bölüm gözden geçirildi.
\item \textbf{Kalkınma Planı ve 2030 Stratejisi} kaynak PDF'lerden okunup madde numaraları doğrulandı (270, 271; ``Fizik ve Astronomi'').
\item \textbf{Ön çalışma daraltıldı}: grafikler ve ölçüm sonuçları formdan çıkarıldı; yalnızca Tablo 1 kaldı (senin kararın).
\item \textbf{Ekip}: Bahar Cayib proje ortağı olarak eklendi; görevler takvimde isimle yazıldı.
\item \textbf{Son kontrol (7 Kasım)}: 35 DOI Crossref'te, kaynaklara yüklenen iddialar makale özetlerinde, Wang vd.'nin Tablo 3'ü ve ölçütleri tam metinde, veri arşivleri tek tek kontrol edildi; Tablo 1 arşivlerden baştan üretildi. Bulunan ve düzeltilenler:
  \begin{itemize}
  \item KELT-16 b'nin keşif makalesinde RV olduğu bulundu $\rightarrow$ \sR{} = 39,3 hesaplandı; ``veri yok'' yerine ``veri yetersiz''. Sonuç değişmedi.
  \item Baştürk vd. 2025 arXiv'den MNRAS yayınına güncellendi; Winn vd. 2009'un dergi adı düzeltildi; Oberst vd. 2017 eklendi (35 kaynak).
  \item ``Zaman tabanı 2--5 yıl'' $\rightarrow$ ``2--6 yıl'' (WASP-4 için 5,9 yıl).
  \end{itemize}
\end{enumerate}

\subsection{Verilen kararlar ve gerekçeleri}
\begin{longtable}{@{}p{4.2cm}p{11.5cm}@{}}
\toprule
Karar & Gerekçe \\ \midrule\endhead
Yalnızca arşiv RV'si & Tayfçeker yok; yayımlanmış RV'ler yeterli (Bölüm 1). \\
Bütçe 0 TL & Tüm veri ve yazılım ücretsiz; uydurma kalem tutarlılığı bozar. Bütçe değerlendirme kalemi değil. \\
LOOCV kapsam dışı & Hocanın direktifi; LOOCV yalnızca Wang'ın listesi olarak kullanıldı. \\
Dört hedef & Ölçütlerin sonucu; her biri farklı bir durumu temsil ediyor. \\
Ön çalışma = yalnız Tablo 1 & Savunabileceğin ve kendin anlattığın şeyle sınırlı tutmak. \\
Çıktı: kısa makale & Gerçekçi; ``yeterli görülmesi halinde'' koşuluyla. \\
Marka/kurum adı yok & Çağrı madde 4.1.8. Arşiv adları zorunlu veri kaynağı olduğu için var. \\
Bahar proje ortağı & Koşulları sağlıyor; destek görevleri, zor kısımlar sende. \\ \bottomrule
\end{longtable}

\subsection{Dosyalar nerede? (rraeyz/proje deposu)}
\begin{itemize}
\item \texttt{2209-A\_Arastirma\_Onerisi\_Formu\_Idris\_Avlar.docx}: hocaya gidecek Word formu.
\item \texttt{TUBITAK\_2209A\_form\_metinleri.md}: aynı metnin düz hali (TYBS'ye buradan da kopyalanabilir).
\item \texttt{kilavuz.pdf}: bu belge.
\item \texttt{kod/secim\_tablosu.py}: Tablo 1'in ilk (hesaplı, \sR) sürümünü üretir; formda artık kullanılmıyor. \texttt{kod/oc\_diyagrami.py}: \OC{} analizi. \texttt{kod/tess\_zaman.py}: TESS zaman ölçümü. \texttt{kod/rv\_analiz.py}: RV eğilimi ve Rømer.
\item \texttt{veri/tablo1\_secim.csv}: Tablo 1'in ham hali ve her değerin kaynağı. \texttt{veri/tess\_*.csv}: denemede ölçülen TESS zamanları. \texttt{veri/rv\_*}: indirilen RV'ler.
\item \texttt{sekiller/}, \texttt{on\_calisma\_colab.ipynb}: eski denemelerden; formda kullanılmıyor.
\end{itemize}

% =====================================================================
\section{Formda olmayan ama yapılmış denemeler}
% =====================================================================
\begin{dikkat}[Bunu nasıl kullanmalısın?]
Bunlar yöntemin çalıştığını görmek için yapıldı, sonra senin kararınla formdan çıkarıldı. Kodları yapay zekâ (Claude) yazdı ve bu oturumda çalıştırdı; sen çalıştırmadın. Formda yazmadığı için kimse senden beklemiyor. Hoca ``yöntemi denedin mi?'' derse: \emph{``Evet, kodları WASP-12 b üzerinde denedim, literatürle uyumlu çıktı; asıl analizi projede yapacağım''} diyebilirsin — ama \textbf{yalnızca aşağıdakini anladıysan}. Anlamadığın sayıyı söyleme.
\end{dikkat}

\subsection{Arşiv RV'lerinden \sR{} hesabı}
Tablo 1'in ilk sürümünde Rømer sütunu, yayımlanmış RV ölçümlerine (Knutson, Bouma, Bonomo, Oberst, DACE) kodla model uydurularak hesaplanmıştı: geçiş sırasındaki ölçümler atıldı, her alete ayrı sıfır noktası verildi, gezegen sinyali + doğrusal eğilim uyduruldu, eğimin belirsizliği $\dot P=P(dv/dt)/c$ ile dönem değişimine çevrildi. Sonuç (\msyil): WASP-12 1,6; WASP-4 0,4; HAT-P-7 1,2; TrES-1 2,2 (seçildi); KELT-16 39,3; TrES-5 29,2; WASP-46 56,3 (elendi). Yani literatüre dayalı seçimle \emph{aynı dört hedef}. Formda artık yok.

\subsection{Yöntem doğrulaması (TESS sektör 20)}
Wang vd.'nin zaten ölçtüğü sektör 20'deki WASP-12 b geçişleri yeniden ölçüldü: 21 geçiş, 20'si ortak. Ortalama fark $-1{,}7$ s, farkların hataya oranının ortancası 0,40σ, ortanca hata 45 s. Yani sistematik kayma yok; ölçüm yöntemi Wang vd. ile uyumlu.

\subsection{Wang sonrası 174 yeni TESS zamanı}
\begin{center}
\begin{tabular}{@{}llrr@{}}
\toprule
Hedef & Sektörler & Yeni zaman & Ortanca hata \\ \midrule
WASP-12 b & 71, 72 & 38 & 43 s \\
WASP-4 b & 36, 69, 96, 103, 105, 106 & 75 & 33 s \\
HAT-P-7 b & 74, 75, 78, 81, 82, 86 & 40 & 56 s \\
TrES-1 b & 74, 78, 80, 81, 86 & 21 & 31 s \\ \midrule
Toplam & & 174 & \\ \bottomrule
\end{tabular}
\end{center}

\subsection{Deneme \OC{} sonuçları}
\begin{center}\small
\begin{tabular}{@{}lrrrl@{}}
\toprule
Hedef & Zaman sayısı & dP/dt (\msyil) & ΔBIC & Rømer karşılaştırması (arşiv RV) \\ \midrule
WASP-12 b & 227 & $-30{,}90\pm0{,}90$ & 1177 & Rømer $\approx -3{,}6\pm1{,}6$ $\rightarrow$ açıklayamaz ($\sim$15σ) \\
WASP-4 b & 197 & $-4{,}79\pm0{,}40$ & 138 & Rømer $\approx -5{,}9$ $\rightarrow$ büyük ölçüde açıklıyor \\
HAT-P-7 b & 117 & $+30{,}00\pm4{,}73$ & 35,5 & Rømer $\approx +13$ $\rightarrow$ kısmen \\
TrES-1 b & 94 & $-4{,}26\pm2{,}63$ & $-1{,}9$ & anlamlı değişim yok \\ \bottomrule
\end{tabular}
\end{center}
Bunlar ön deneme; projede zaman sistemleri tek tek kontrol edilip, tüm testlerle yeniden yapılacak. İlginç olan: dört sonucun dördü de hipotezlerle (H1--H3) aynı yönde. Ama \textbf{formda bu sayılar yok}; hipotezler literatüre dayandırıldı.

% =====================================================================
\section{Kavramlar sözlüğü}
% =====================================================================
\begin{description}[style=nextline,leftmargin=1.2em,itemsep=3pt]
\item[Geçiş (transit)] Gezegenin yıldızın önünden geçmesi; yıldızın ışığında küçük bir düşüş.
\item[Orta geçiş zamanı ($T_c$)] Geçişin tam ortasının zamanı.
\item[Epok ($E$)] Referans geçişten bu yana kaçıncı geçiş olduğu (tam sayı).
\item[Efemeris] Geçiş zamanlarını öngören formül: $T = T_0 + P E$.
\item[\OC] Gözlenen eksi hesaplanan zaman. Dönem sabitse düz, sabit hızla değişiyorsa parabol.
\item[dP/dt, \msyil] Dönemin yılda kaç milisaniye değiştiği. Negatif = kısalma.
\item[BJD\_TDB] Güneş Sistemi kütle merkezine taşınmış, düzgün akan zaman ölçeği. HJD/BJD\_UTC'den farkı dakikalara varabilir.
\item[σ, 3σ] Belirsizlik; bir değer sıfırdan 3σ uzaksa ``anlamlı'' sayılır.
\item[$\chi^2$, indirgenmiş $\chi^2$] Modelle verinin uyumsuzluğu; indirgenmiş $\chi^2\approx1$ hataların gerçekçi olduğunu gösterir.
\item[BIC, ΔBIC] Model seçimi ölçütü; ΔBIC $>10$: çok güçlü tercih.
\item[LOOCV] Birini dışarıda bırak çapraz doğrulama: her seferinde bir ölçümü çıkarıp sonucun dayanıklılığını sınamak.
\item[Gelgit bozulması] Gelgit etkileşimiyle yörüngenin gerçekten küçülmesi.
\item[$Q'_*$] Yıldızın gelgit kalite faktörü; küçükse bozulma hızlı.
\item[Rømer etkisi (ışık-yolu etkisi, LTTE)] Sistemin ivmelenmesinden kaynaklanan görünür dönem değişimi; yörünge değişmez.
\item[Radyal hız (RV)] Yıldızın bize yaklaşma/uzaklaşma hızı (tayftaki Doppler kaymasından).
\item[Rømer öngörüsü] RV ivmesinden beklenen görünür dönem değişimi, $\dot P = P(dv/dt)/c$ (\msyil). Bouma vd. 2020 Tablo 5'te 51 gezegen için yayımlanmış.
\item[Jitter] Yıldızın kendi yüzey hareketlerinden gelen ek RV saçılması.
\item[Rossiter--McLaughlin] Geçiş sırasında RV'deki bozulma.
\item[Apsis hareketi] Dışmerkez yörüngenin uzun ekseninin dönmesi; \OC'de sinüs.
\item[Applegate etkisi] Yıldızın manyetik döngüsünün dönemi çevrimsel değiştirmesi.
\item[Üçüncü cisim] Sistemdeki ek bir gezegen/yıldız; Rømer ya da dinamik etki yaratır.
\item[Roche sınırı] Gezegenin gelgitle parçalanmaya başladığı uzaklık.
\item[TESS, sektör, SPOC] NASA'nın geçiş arayan uydusu; bir sektör $\sim$27 gün; SPOC: 2 dakikalık işlenmiş ışık eğrileri.
\item[batman] Geçiş ışık eğrisi modelleyen Python paketi.
\item[MCMC, emcee] Parametre belirsizliklerini örnekleme ile bulan yöntem ve paketi.
\item[Lomb--Scargle] Düzensiz aralıklı veride periyot arama yöntemi.
\item[VizieR, MAST, DACE, NASA Exoplanet Archive, ETD/VarAstro, ExoClock] Kullanılan açık arşivler (katalog tabloları, TESS verisi, RV, sistem parametreleri, yer gözlemi zamanları).
\end{description}

% =====================================================================
\section{Son kontrol listesi}
% =====================================================================
\begin{itemize}[label=$\square$]
\item ARBİS: benim ve Bahar'ın kaydı güncel.
\item Hocanın YÖKSİS/ARBİS bilgileri güncel.
\item Hoca Word belgesini gördü ve onayladı.
\item TYBS: başlık, danışman, proje ortağı girildi.
\item Bütçe: 0 TL.
\item Metin kutuları dolduruldu; tablolar ``Tablo'' düğmesiyle kuruldu; özel karakterler doğru görünüyor.
\item Kaynakça: 35 kaynak, DOI'ler var.
\item Yapay zekâ beyanı (kendim).
\item Önizleme kontrol edildi, hiçbir bölüm boş değil.
\item Onaylandı $\rightarrow$ danışman onayı $\rightarrow$ kuruluş e-imzası (11 Kasım 17.30'dan önce).
\item Bölüm 1'i ve Bölüm 3'teki soruları bir kez sesli anlattım.
\end{itemize}

% =====================================================================
\appendix
\section{TYBS'ye yapıştırılacak form metni (tamamı)}
Bu ek, \texttt{TUBITAK\_2209A\_form\_metinleri.md} dosyasından otomatik üretilir; Word formuyla aynı metindir.
\subsection*{A. Genel Bilgiler}
{\footnotesize\begin{longtable}{|p{0.450\linewidth}|p{0.450\linewidth}|}\hline
\textbf{Alan} & \textbf{Değer}\\\hline\endhead
Başvuru Sahibinin Adı Soyadı & İdris Avlar\\\hline
Araştırma Önerisinin Başlığı & Sıcak Jüpiterlerde Yörünge Dönemi Değişimlerinin O−C Analiziyle İncelenmesi ve Fiziksel Kökenlerinin Belirlenmesi\\\hline
Danışmanın Adı Soyadı & Hüseyin Er\\\hline
Araştırmanın Yürütüleceği Kurum & Atatürk Üniversitesi, Fen Fakültesi, Astronomi ve Uzay Bilimleri Bölümü\\\hline
\end{longtable}}
\subsection*{Özet}
Yıldızlarına çok yakın yörüngelerde dolanan sıcak Jüpiterlerin yörünge dönemlerinin zamanla değişip değişmediği, yıldız–gezegen gelgit etkileşimini doğrudan sınayabilen az sayıdaki gözlemsel yoldan biridir. Wang vd. (2024), TESS verisiyle 326 sıcak Jüpiteri inceleyerek uzun dönemli dönem değişimi gösteren 26 aday bildirmiş; bunlardan 14'ü, tek bir ölçümün çıkarılmasına karşı sonucun dayanıklılığını sınayan “birini dışarıda bırak çapraz doğrulama” (LOOCV) testini geçmiştir. Ancak geçiş zamanlarında görülen parabolik bir değişim yalnızca gelgit kaynaklı yörünge bozulmasından değil; apsis hareketi, uzak bir üçüncü cismin ışık-yolu (Rømer) etkisi ya da kütle çekimsel etkisi, Applegate mekanizması, kütle kaybı ve veri hatalarından da kaynaklanabilir. Bu projede LOOCV testini geçen 14 aday, TESS veri yeterliliği ve fiziksel kökenin ayırt edilebilmesi için gerekli arşiv radyal hız verisinin yeterliliği ölçütleriyle elenmiş; WASP-12 b, WASP-4 b, HAT-P-7 b ve TrES-1 b hedef olarak seçilmiştir. Her hedef için literatür ve katalog geçiş zamanları derlenecek, tümü BJD\_TDB zaman sistemine dönüştürülecek ve Wang vd. tarafından kullanılmamış yeni TESS sektörlerinden orta geçiş zamanları ölçülecektir. Gözlenen ve doğrusal efemerisle hesaplanan zamanların farkı olan O−C diyagramlarına doğrusal, kuadratik ve kübik modeller uydurulacak; dönem değişiminin varlığı, değişim hızının anlamlılığı ve Bayes Bilgi Ölçütü ile sınanacaktır. Anlamlı değişim bulunan hedeflerde, yalnızca zamanlamayla ayırt edilemeyen gelgit ve Rømer etkileri, yayımlanmış arşiv radyal hız ölçümlerinden öngörülen dönem değişimiyle karşılaştırılarak ayrılacak; diğer etkiler gelgit kalite faktörü, örtülme zamanları, dışmerkezlik, yıldız etkinliği, periyodogram ve Roche sınırı testleriyle değerlendirilecektir. Projede yeni gözlem yapılmayacak; tüm veriler açık arşivlerden sağlanacaktır. Ön çalışmada 14 adayın TESS ve radyal hız verilerinin yeterliliği incelenerek hedef seçimi yapılmıştır. Bütçe gerektirmeyen 12 aylık projenin çıktıları dört hedef için güncel efemerisler, fiziksel köken hükümleri ve açık erişimli veri seti ile analiz kodudur; sonuçların danışmanla birlikte kısa bir makale olarak paylaşılması hedeflenmektedir.

\textbf{Anahtar Kelimeler:} Sıcak Jüpiter, geçiş zamanlaması, O−C diyagramı, yörünge bozulması, TESS

\subsection*{1. Araştırma Önerisinin Bilimsel Niteliği}
\subsubsection*{1.1. Konunun Önemi ve Araştırma Önerisinin Bilimsel Niteliği}
Sıcak Jüpiterler, kütleleri Jüpiter mertebesinde olan ve yıldızlarının etrafında birkaç günde dolanan gezegenlerdir. Yıldızlarına çok yakın oldukları için yıldızda gelgit kabarıklıkları oluştururlar. Yıldızın kendi ekseni etrafındaki dönmesi yörünge hareketinden yavaş olduğunda bu kabarıklık gezegeni geriye doğru çeker; gezegen açısal momentum kaybederek içeri doğru sarmallanır ve yörünge dönemi kısalır. Bu sürecin hızı, yıldızın gelgit enerjisini ne kadar verimli soğurduğunu gösteren değiştirilmiş gelgit kalite faktörüne (Q′*) bağlıdır. Kuramsal ve görgül tahminler Q′* için birkaç mertebelik bir aralık verdiğinden (Penev vd., 2018), dönem değişiminin doğrudan ölçümü yıldız içi fiziği sınamanın ender yollarından biridir.

Dönem değişimi, geçiş zamanlarının gözlenen (O) ve sabit dönemle hesaplanan (C) değerleri arasındaki farkın zamana karşı çizildiği O−C diyagramıyla incelenir. Dönem sabit bir hızla değişiyorsa O−C bir parabol çizer. Bugüne kadar ikna edici biçimde doğrulanmış tek örnek WASP-12 b'dir; geçiş, örtülme ve radyal hız verileri birlikte kullanılarak apsis hareketi ve ivme açıklamaları dışlanmış, dönem değişim hızı yaklaşık −29 ms/yıl ve Q′* ≈ 1,8×10⁵ bulunmuştur (Patra vd., 2017; Yee vd., 2020; Winn ve Stefánsson, 2025). Wang vd. (2024), 326 sıcak Jüpiterin TESS (Ricker vd., 2015) geçişlerini literatür zamanlarıyla birleştirerek 26 dönem değişimi adayı bildirmiş, bunlardan 12'sinin sonucunun tek bir ölçümün çıkarılmasıyla kaybolduğunu göstermiş ve gelgit dışında apsis hareketi, Rømer etkisi ve yıldızın manyetik etkinlik döngüsüne bağlı Applegate etkisinin (Applegate, 1992; Watson ve Marsh, 2010) de katkı verebileceğini belirtmiş; ancak adayların fiziksel kökenini sistem sistem sınamamıştır.

Literatürdeki temel sorun, O−C diyagramındaki parabolün tek başına kökenini göstermemesidir. Uzak bir eşlikçi nedeniyle bakış doğrultusunda sabit ivmelenen bir sistemde ışığın yol süresi değişir ve geçiş zamanları gelgit bozulmasıyla matematiksel olarak aynı kuadratik izi bırakır (Bouma vd., 2020). Nitekim WASP-4 b için aynı dönem değişimi bir çalışmada gelgit bozulmasıyla (Baştürk vd., 2025), başka bir çalışmada yaklaşık 7 Jüpiter kütleli uzak bir eşlikçinin ışık-yolu etkisiyle açıklanmıştır (Winn ve Stefánsson, 2025). HAT-P-7 b'de radyal hızlarda uzun dönemli bir eğilim bulunmuş (Winn vd., 2009) ve geçiş zamanlaması ile radyal hızların birlikte kullanılmasıyla yaklaşık 32 AB uzaklıkta eksantrik yörüngeli bir yıldız eşlikçi önerilmiştir (Yang vd., 2025). TrES-1 b için ise farklı çalışmalar birbiriyle uyumsuz dönem değişim hızları ve yorumlar ortaya koymuştur (Adams vd., 2024; Hagey vd., 2025; Kutluay vd., 2026). Ayrıca literatür zamanlarındaki hatalar, zaman sistemi karışıklıkları, gerçekçi olmayan küçük hata çubukları ve tek tek aykırı ölçümler sahte sinyaller üretebilmektedir (Adams vd., 2024; Rothmeier vd., 2025; Sodickson ve Grunblatt, 2025). Bu nedenle adayların güncel veriyle yeniden sınanması ve olası her fiziksel etkinin aynı yöntemle sistematik biçimde değerlendirilmesi gerekmektedir.

Bu proje, Wang vd. (2024)'nin LOOCV testini geçen 14 adayından önceden tanımlanmış ölçütlerle seçilen dört sıcak Jüpiterde (WASP-12 b, WASP-4 b, HAT-P-7 b, TrES-1 b) bu boşluğu doldurmayı amaçlamaktadır. Hedefler, TESS veri yeterliliği ve arşiv radyal hız verisinin yeterliliği ölçütleriyle seçilmiştir (bkz. Yöntem, Tablo 1). Seçilen hedefler birbirini tamamlayan dört durumu temsil etmektedir: WASP-12 b doğrulanmış gelgit bozulmasıyla yöntemin kalibrasyonunu, WASP-4 b gelgit ile ışık-yolu etkisinin tartışıldığı bir vakayı, HAT-P-7 b gelgitle açıklanamayan bir dönem artışını, TrES-1 b ise bildirilen dönem kısalmasının literatür zamanlarının güvenilirliğine bağlı olup olmadığının sınandığı bir vakayı. Projenin özgün yanları şunlardır: (i) Wang vd. sonrasında yayımlanan TESS sektörlerinden ölçülecek yeni geçiş zamanlarıyla zaman tabanı yaklaşık 2–6 yıl uzatılacaktır; (ii) gelgit etkileşimi, apsis hareketi, Rømer etkisi, üçüncü cismin dinamik etkisi, Applegate etkisi ve kütle kaybı her hedef için önceden yazılmış bir karar ağacıyla değerlendirilecektir; (iii) gelgit ile Rømer etkisi, farklı arşivlerde yayımlanmış radyal hız ölçümlerinden öngörülen dönem değişimiyle karşılaştırılarak ayrılacaktır.

\textbf{Araştırma sorusu:} Wang vd. (2024)'nin LOOCV testini geçen adaylardan seçilen dört sıcak Jüpiterin yörünge dönemleri güncel TESS verisiyle anlamlı biçimde değişmekte midir; değişiyorsa bu değişimden gelgit etkileşimi, apsis hareketi, Rømer etkisi, üçüncü cisim, Applegate etkisi ve kütle kaybından hangisi sorumludur?

\textbf{Hipotezler:}

H1: WASP-12 b'de güncel TESS verisiyle bulunan dönem değişim hızı literatür değeriyle (Yee vd., 2020) 2σ içinde uyumludur ve arşiv radyal hızlarından öngörülen Rømer katkısı bu değişimi açıklayamaz (gelgit kökenli bozulma).

H2: WASP-4 b ve HAT-P-7 b'deki dönem değişimlerinin önemli bir kısmı, arşiv radyal hız ölçümlerinin gösterdiği ivmeden kaynaklanan ışık-yolu (Rømer) etkisiyle açıklanır.

H3: TrES-1 b'de bildirilen dönem kısalması, literatür zamanlarının belirsizlikleri gerçekçi biçimde değerlendirildiğinde ve yeni TESS zamanları eklendiğinde istatistiksel olarak anlamlı değildir.

Proje, On İkinci Kalkınma Planı'nın (2024–2028) temel bilimlere özel önem verilmesini ve uzay araştırmalarını kritik alanlar arasında sayan hedefiyle (madde 270) ve genç nesillerin uzay alanında çalışmaya teşvik edilerek uzay araştırmalarının yaygınlaştırılmasını öngören hedefiyle (madde 271) doğrudan ilişkilidir (T.C. Cumhurbaşkanlığı Strateji ve Bütçe Başkanlığı, 2023). 2030 Sanayi ve Teknoloji Stratejisi de TÜBİTAK araştırma desteklerinin astronomi dahil temel bilimlerde bilgi üretimini artıracak şekilde yaygınlaştırılmasını öngörmektedir (T.C. Sanayi ve Teknoloji Bakanlığı, 2025). Proje, uzay görevi verilerinin (TESS) işlenmesi ve yorumlanması konusunda lisans düzeyinde araştırmacı yetiştirilmesine katkı sağlayarak bu hedeflere hizmet etmektedir.

\subsubsection*{1.2. Amaç ve Hedefler}
Projenin amacı, Wang vd. (2024)'nin LOOCV testini geçen 14 sıcak Jüpiter adayından önceden tanımlanmış ölçütlerle seçilen dört hedefte yörünge dönemi değişiminin varlığını güncel TESS ve katalog verileriyle sınamak ve değişim varsa fiziksel kökenini belirlemektir.

\textbf{Hedef 1 (Veri seti):} Dört hedefin her biri için en az 70 geçiş zamanından oluşan, tümü BJD\_TDB zaman sistemine dönüştürülmüş ve kaynakları etiketlenmiş bir zaman tablosu oluşturmak; Wang vd. sonrası TESS sektörlerinden toplam en az 150 yeni orta geçiş zamanını, ortanca belirsizliği 60 saniyeyi aşmayacak biçimde ölçmek.

\textbf{Hedef 2 (Dönem değişimi):} Her hedef için O−C diyagramını çizmek; doğrusal, kuadratik ve kübik modelleri karşılaştırarak dönem değişim hızını (dP/dt) belirsizliğiyle ve ΔBIC değeriyle raporlamak ve her hedef için “anlamlı değişim var/yok” kararını vermek. Yöntemin doğruluğu, WASP-12 b için literatür değeriyle 2σ içinde uyumla gösterilecektir.

\textbf{Hedef 3 (Fiziksel köken):} Anlamlı değişim bulunan her hedef için altı fiziksel etkinin her birine (gelgit, apsis hareketi, Rømer, üçüncü cisim, Applegate, kütle kaybı) önceden tanımlanmış ölçütlerle “dışlandı / olası değil / olası / en olası / ayrılamadı” hükmü veren gerekçeli bir karar tablosu oluşturmak.

\textbf{Hedef 4 (Efemeris):} Dört hedef için güncel efemerisleri (T0, P ve varsa dP/dt) yayımlamak.

\subsection*{2. Yöntem}
\textbf{Araştırma tasarımı.} Proje, açık arşiv verisine dayanan, hipotez sınayan gözlemsel bir çalışmadır. Yeni teleskop ya da tayf gözlemi yapılmayacak; tüm veriler TESS ışık eğrileri için MAST, katalog tabloları için VizieR (Ochsenbein vd., 2000), sistem parametreleri için NASA Exoplanet Archive (Akeson vd., 2013), radyal hızlar için DACE ve yayımlanmış makale tabloları, yer gözlemi geçiş zamanları için ETD/VarAstro (Poddaný vd., 2010) ve ExoClock (Kokori vd., 2023) veritabanlarından sağlanacaktır. Bağımsız değişken geçiş epoku (E), bağımlı değişken orta geçiş zamanı ve ondan türetilen O−C değeridir; türetilen temel büyüklük dönem değişim hızıdır (dP/dt). Kontrol edilen değişkenler zaman sistemi, ölçüm belirsizliği ve veri kaynağıdır.

\textbf{Örneklem seçimi.} Hedefler, bu projenin kendi O−C sonuçlarından bağımsız üç ölçütle seçilmiştir (Tablo 1). (i) LOOCV ölçütü: Wang vd. (2024)'nin LOOCV testini geçmiş olmak (14 aday). (ii) TESS verisi ölçütü: Wang vd. tarafından raporlanan TESS geçiş zamanı sayısının en az 30 olması; adaylar arasındaki en kısa yörünge döneminde bile bu sayı en az iki TESS sektörüne yayılan veri anlamına gelir (7 aday). (iii) Radyal hız verisi ölçütü: Fiziksel kökenin belirlenebilmesi için, yayımlanmış arşiv radyal hız ölçümlerinden öngörülen Rømer katkısının belirsizliğinin Wang vd.'nin raporladığı |dP/dt| değerinin üçte birinden küçük olması (4 aday: WASP-12 b, WASP-4 b, HAT-P-7 b, TrES-1 b).

\textbf{İP1 – Zaman tablosunun oluşturulması (Ay 1–2).} Geçiş zamanları Ivshina ve Winn (2022) veritabanından, Wang vd. (2024)'nin TESS ölçümlerinden, ETD/VarAstro (Poddaný vd., 2010) ve ExoClock (Kokori vd., 2023) veritabanlarından ve hedefe özel makalelerden (Yee vd., 2020; Bouma vd., 2020; Baştürk vd., 2025; Hagey vd., 2025) derlenecektir. Her zamanın sistemi makalesinden kontrol edilecek, HJD ve BJD\_UTC zamanları Astropy (Astropy Collaboration vd., 2022) ile BJD\_TDB'ye dönüştürülecektir. Zaman sistemi makalesinden belirlenemeyen zamanlar işaretlenecek; analiz bu zamanlar dahil ve hariç olmak üzere iki kez yapılarak sonucun bu zamanlara duyarlılığı raporlanacaktır. ETD verilerinde yalnızca yüksek kalite notlu tam geçişler kullanılacaktır. Epok numarasının kesirli kısmı 0,02'yi aşan zamanlar ayrıca kontrol edilecektir.

\textbf{İP2 – Yeni TESS zamanlarının ölçülmesi (Ay 3–4).} Wang vd. sonrası SPOC 2 dakikalık ışık eğrileri MAST'tan indirilecek, kalite bayraklı noktalar atılarak normalize edilecektir. Tüm geçişler katlanarak geçiş şekli (Rp/R*, a/R*, eğiklik) batman modeliyle (Kreidberg, 2015) belirlenecek; ardından her geçiş, yalnızca orta zaman ve yerel taban çizgisi serbest bırakılarak uydurulacak ve orta zamanın belirsizliği emcee ile MCMC örneklemesinden (Foreman-Mackey vd., 2013) elde edilecektir. Kenar kararması katsayıları sabit tutulacaktır; geçiş biçimi simetrik olduğundan orta zaman ölçümüne etkisi ihmal edilebilir. Kısmi geçişler kullanılmayacaktır. Yöntem, Wang vd.'nin ölçtüğü sektörlerden biri yeniden ölçülerek ortak geçişlerde karşılaştırılacak ve doğrulanacaktır.

\textbf{İP3 – O−C analizi ve model seçimi (Ay 5–6).} Doğrusal (T = T0 + P·E), kuadratik (T = T0 + P·E + ½·(dP/dE)·E²) ve kübik efemerisler ağırlıklı en küçük karelerle uydurulacak, dP/dt = (1/P)(dP/dE) hesaplanacaktır. Aykırı noktalar 5σ yinelemeli kırpmayla belirlenecektir. Literatür ve TESS zamanlarının hata çubukları, her grubun kuadratik model artıklarından hesaplanan indirgenmiş χ² değeri 1 olacak şekilde (yalnızca büyütülerek) ölçeklenecektir. Model seçimi Bayes Bilgi Ölçütü (BIC = χ² + k·ln n; Schwarz, 1978) ile yapılacak; Wang vd. (2024) ile aynı ölçütler (|dP/dt|/σ ≥ 3 ve ΔBIC > 10) uygulanacaktır. Kübik model, dönem değişim hızının zamanla sabit olup olmadığını sınamak için kullanılacaktır.

\textbf{İP4 – Fiziksel kökenin belirlenmesi (Ay 7–11).} (a) Rømer etkisi: Yayımlanmış arşiv radyal hızlarına (Knutson vd., 2014; Bonomo vd., 2017; Bouma vd., 2020; Yee vd., 2020; DACE) her alet için ayrı sıfır noktası, gezegenin dairesel yörünge sinyali ve doğrusal eğilim uydurulacaktır. Geçiş sırasında alınan ölçümler (Rossiter–McLaughlin etkisi) çıkarılacak, hatalar yıldız kaynaklı saçılmayı içerecek şekilde ölçeklenecek, doğrusal eğilim yetersiz kalırsa ikinci dereceden terim eklenecektir. Rømer etkisinin öngördüğü dönem değişimi dP/dt = P·(dv/dt)/c ile hesaplanıp O−C sonucuyla karşılaştırılacaktır (fark < 2σ: Rømer açıklayabilir; 2–3σ: kısmen; > 3σ: Rømer tek başına açıklayamaz). Sonucun kullanılan radyal hız veri setine duyarlılığı ayrıca raporlanacaktır. (b) Gelgit: NASA Exoplanet Archive parametreleriyle dP/dt = −(27π/2Q′*)(Mp/M*)(R*/a)⁵ bağıntısından (Patra vd., 2017) gereken Q′* hesaplanacak ve kuramsal aralıkla karşılaştırılacaktır; yıldızın yörüngeden yavaş döndüğü sistemlerde gelgit dönemi yalnızca kısaltabileceğinden, dönemi uzayan hedefte gelgit dışlanır (Wang vd., 2024). (c) Apsis hareketi: dışmerkezlik üst sınırından beklenen genlik (≈ eP/π) hesaplanacak; örtülme zamanlarının bulunduğu WASP-12 b'de (Yee vd., 2020) örtülmelerin geçişlere göre kayma yönü sınanacaktır. (d) Applegate etkisi ve üçüncü cismin dinamik etkisi: kuadratik model artıklarında Lomb–Scargle periyodogramıyla (Lomb, 1976; Scargle, 1982; Astropy Collaboration vd., 2022) periyodik sinyal aranacak (yanlış alarm olasılığı < 0,01; gözlem mevsiminden kaynaklanan yaklaşık bir yıllık tepeler sinyal sayılmayacak), kübik modelle eğriliğin sabitliği sınanacak, yıldız etkinliği literatürden değerlendirilecektir. (e) Kütle kaybı: a/a\_Roche oranı hesaplanacaktır. Her etki için hüküm önceden yazılmış karar ağacıyla verilecektir.

\textbf{Ön çalışma.} Ön çalışmada Wang vd. (2024)'nin LOOCV testini geçen 14 adayın, projenin yürütülebilmesi için gereken verilerinin yeterliliği incelenmiştir. Her aday için Wang vd.'nin TESS geçiş zamanı sayısı belirlenmiş; yayımlanmış radyal hız ölçümleri (Knutson vd., 2014; Bonomo vd., 2017; Oberst vd., 2017; Bouma vd., 2020; DACE) derlenerek Rømer katkısının hangi duyarlılıkla öngörülebileceği hesaplanmıştır. Bu inceleme sonucunda seçim ölçütlerini sağlayan dört hedef belirlenmiş (Tablo 1) ve dört hedefin geçiş zamanlarının ve Wang vd. sonrası TESS ışık eğrilerinin açık arşivlerde erişilebilir olduğu doğrulanmıştır. Yöntem adımları çalışma takvimindeki İP1–İP4 ile birebir ilişkilidir.

\subsection*{3. Proje Yönetimi}
\subsubsection*{3.1. Çalışma Takvimi}
{\footnotesize\begin{longtable}{|p{0.207\linewidth}|p{0.207\linewidth}|p{0.207\linewidth}|p{0.207\linewidth}|}\hline
\textbf{Tarih Aralığı} & \textbf{Faaliyetler} & \textbf{Kim(ler)} & \textbf{Başarı Ölçütü ve Katkısı}\\\hline\endhead
1.–2. ay & İP1: Arşiv ve literatür geçiş zamanı veri setinin oluşturulması; tüm zamanların BJD\_TDB'ye dönüştürülmesi ve kaynak etiketlemesi & İdris Avlar (proje yürütücüsü; zaman sistemi dönüşümleri); Bahar Cayib (proje ortağı; literatür ve arşiv zamanlarının derlenmesi); Hüseyin Er (akademik danışman) & Her hedef için en az 70 zaman; tüm zamanlar BJD\_TDB'de; zaman sistemi belirsiz kaynaklar ayrıca işaretli (\%20)\\\hline
3.–4. ay & İP2: Wang vd. sonrası TESS sektörlerinden orta geçiş zamanlarının ölçülmesi ve Wang vd.'nin ölçtüğü bir sektörün yeniden ölçülerek yöntemin doğrulanması & İdris Avlar (proje yürütücüsü; ışık eğrisi modellemesi ve zaman ölçümü); Bahar Cayib (proje ortağı; TESS verilerinin indirilmesi ve geçişlerin görsel kontrolü) & Dört hedef için en az 150 yeni zaman; ortanca belirsizlik en fazla 60 s; Wang vd. ile ortak geçişlerde ortanca fark en fazla 1σ (\%20)\\\hline
5.–6. ay & İP3: O−C diyagramları; hata çubuklarının ölçeklenmesi; doğrusal, kuadratik ve kübik modeller; BIC ile model seçimi & İdris Avlar (proje yürütücüsü; model uydurma ve istatistiksel karşılaştırma); Bahar Cayib (proje ortağı; O−C diyagramlarının hazırlanması ve veri kontrolü); Hüseyin Er (akademik danışman) & Dört hedef için dP/dt ± σ, ΔBIC ve “anlamlı değişim var/yok” kararı; WASP-12 b'de literatürle 2σ uyum (\%25)\\\hline
7.–11. ay & İP4: Arşiv radyal hızlarıyla Rømer testi; gelgit (Q′*), apsis, Applegate, üçüncü cisim ve kütle kaybı testleri; karar tablosu & İdris Avlar (proje yürütücüsü; fiziksel testler ve karar tablosu); Bahar Cayib (proje ortağı; radyal hız verilerinin ve yıldız/gezegen parametrelerinin derlenmesi); Hüseyin Er (akademik danışman) & Anlamlı değişim bulunan her hedef için altı etkinin her birine gerekçeli hüküm verilmiş karar tablosu (\%35)\\\hline
\end{longtable}}
\subsubsection*{3.2. Risk Yönetimi}
{\footnotesize\begin{longtable}{|p{0.450\linewidth}|p{0.450\linewidth}|}\hline
\textbf{En Önemli Riskler} & \textbf{Alınacak Tedbirler (B Planı)}\\\hline\endhead
Literatür zamanlarında zaman sistemi belirsizliği veya gerçekçi olmayan küçük hata çubukları sahte sinyal üretebilir. & Her kaynağın zaman sistemi makalesinden kontrol edilecek; belirlenemeyen zamanlarla ve onlarsız iki ayrı analiz yapılacak; hata çubukları her veri grubu için indirgenmiş χ² ile ölçeklenecektir.\\\hline
Yeni TESS sektörlerinden bazı geçişlerin zamanı ölçülemeyebilir ve nitelikli zaman sayısı hedefin altında kalabilir (veri boşluğu, düşük sinyal-gürültü oranı). & Tek tek ölçülemeyen geçişler için sektördeki tüm geçişler birlikte uydurularak sektör başına tek bir orta geçiş zamanı elde edilecektir; yöntem değişikliği yalnızca bu adımla sınırlıdır, sonraki analizler aynı kalır. En kötü durumda Wang vd. ve Ivshina ve Winn TESS zamanlarıyla devam edilecektir.\\\hline
Bir hedefte güncel veriyle anlamlı dönem değişimi bulunmayabilir. & Bu durum araştırma sorusunun (“değişim var mı?”) bir cevabı olarak raporlanacak ve güncel sabit dönemli efemeris verilecektir; fiziksel köken analizi anlamlı değişim bulunan hedeflerle sürdürülecektir.\\\hline
Farklı arşivlerdeki radyal hız veri setleri, alet sıfır noktası farkları nedeniyle farklı eğimler verebilir. & Tüm yayımlanmış ölçümler alet başına ayrı sıfır noktasıyla birlikte uydurulacak ve sonucun veri setine duyarlılığı raporlanacaktır; yeterli duyarlılık sağlanamazsa literatürde yayımlanmış eğimler ve radyal hız gerektirmeyen testler (Q′*, eğrilik, periyodogram, Roche) kullanılacaktır.\\\hline
Gelgit ve Rømer etkileri mevcut veriyle ayrılamayabilir. & İki senaryonun da büyüklüğü verilecek ve “ayrılamadı” hükmü gerekçesiyle raporlanacaktır.\\\hline
Veri arşivlerine erişim geçici olarak kesilebilir. & Aynı katalogların ayna sunucuları kullanılacak; indirilen tüm veriler yerel olarak saklanacaktır.\\\hline
İş paketlerinde gecikme yaşanabilir. & İş paketleri ardışık planlanmış ve 12. ay tampon olarak bırakılmıştır; ilerleme akademik danışmanla düzenli görüşmelerle izlenecektir.\\\hline
Yazılım veya bilgisayar arızası. & Tüm kod ve veriler sürüm kontrolüyle çevrim içi bir depoda yedeklenecek; analizler ücretsiz bulut tabanlı bir Python ortamında da çalışabilecek şekilde hazırlanacaktır.\\\hline
\end{longtable}}
\subsubsection*{3.3. Araştırma Olanakları}
{\footnotesize\begin{longtable}{|p{0.450\linewidth}|p{0.450\linewidth}|}\hline
\textbf{Kuruluşta Bulunan Altyapı/Ekipman Türü, Modeli} & \textbf{Projede Kullanım Amacı}\\\hline\endhead
Proje yürütücüsünün kişisel dizüstü bilgisayarı ile üniversite kütüphanesindeki bilgisayarlar ve internet erişimi & Veri indirme, ışık eğrisi analizi, O−C ve radyal hız analizleri\\\hline
Açık erişimli veri arşivleri: MAST (TESS), VizieR, NASA Exoplanet Archive, DACE, ETD/VarAstro, ExoClock & Geçiş zamanları, ışık eğrileri, sistem parametreleri ve radyal hız verileri\\\hline
Açık kaynak yazılımlar: Python, NumPy, Astropy, astroquery, batman, emcee & Tüm analizlerin ücretsiz ve yeniden üretilebilir biçimde yapılması\\\hline
\end{longtable}}
\subsection*{4. Araştırma Önerisinin Yaygın Etkisi}
{\footnotesize\begin{longtable}{|p{0.450\linewidth}|p{0.450\linewidth}|}\hline
\textbf{Çıktı, Etki ve Kazanımlar} & \textbf{Öngörülen Çıktı(lar), Etki(ler) ve Kazanım(lar)}\\\hline\endhead
Bilimsel/Akademik Çıktılar & Dört hedefin güncel O−C analizini ve fiziksel köken karar tablosunu içeren proje sonuç raporu. Sonuçların, yeterli görülmesi halinde danışmanla birlikte kısa bir makale olarak sunulması hedeflenmektedir.\\\hline
Ekonomik/Ticari/Sosyal Çıktılar & Dört hedefin tüm geçiş zamanlarını (yeni TESS zamanları dahil) BJD\_TDB sisteminde içeren açık erişimli veri seti, analiz kodu ve dört hedef için güncel efemerisler; bu çıktılar gelecekteki geçiş gözlemlerinin planlanmasında kullanılabilecektir. Danışmanın uygun görmesi halinde bölüm öğrencilerine yönelik kısa bir O−C analizi tanıtım sunumu yapılması hedeflenmektedir.\\\hline
Yeni Proje(ler) Oluşturmasına Yönelik Çıktılar & Aynı yöntemin Wang vd. (2024)'nin diğer adaylarına genişletilmesiyle olası bir lisansüstü tez çalışmasına temel oluşturulması.\\\hline
\end{longtable}}
\subsection*{5. Belirtmek İstediğiniz Diğer Konular}
Aşağıdaki tablo, yöntem bölümünde özetlenen ön çalışmanın sonucudur.

\textbf{Tablo 1.} Bu çalışmada oluşturulan hedef seçim tablosu. Satırlar, Wang vd. (2024)'nin dönem değişimi bildirdiği 26 adaydan, aynı çalışmanın LOOCV testini geçen 14'üdür (Wang vd., 2024, Tablo 3). Sütunların kaynakları: |dP/dt|/σ (Wang vd.'nin bildirdiği dönem değişim hızının belirsizliğine oranı) ve TESS geçiş sayısı Wang vd. (2024)'nin VizieR'de yayımlanan tablolarından (J/ApJS/270/14) alınmıştır. σ\_Rømer (uzak bir eşlikçinin ivmesinden kaynaklanabilecek dönem değişiminin arşiv radyal hız verisiyle öngörülebildiği belirsizlik, ms/yıl) bu çalışmada, yayımlanmış radyal hız ölçümlerinden (Knutson vd., 2014; Bonomo vd., 2017; Oberst vd., 2017; Bouma vd., 2020; DACE) hesaplanmış ve her aday için en hassas veri seti kullanılmıştır. Seçim ölçütleri: (1) TESS geçiş sayısı en az 30; (2) σ\_Rømer, Wang vd.'nin bildirdiği |dP/dt| değerinin üçte birinden (“Eşik” sütunu) küçük. “–”: ilk ölçütte elendiği için hesaplanmadı.

{\footnotesize\begin{longtable}{|p{0.127\linewidth}|p{0.127\linewidth}|p{0.127\linewidth}|p{0.127\linewidth}|p{0.127\linewidth}|p{0.127\linewidth}|}\hline
\textbf{Aday} & \textbf{|dP/dt|/σ} & \textbf{TESS geçiş sayısı} & \textbf{σ\_Rømer (ms/yıl)} & \textbf{Eşik (ms/yıl)} & \textbf{Sonuç}\\\hline\endhead
CoRoT-2 b & 7,3 & 5 & – & – & Elendi: TESS geçiş sayısı 30'dan az\\\hline
HAT-P-37 b & 5,1 & 8 & – & – & Elendi: TESS geçiş sayısı 30'dan az\\\hline
KELT-16 b & 5,8 & 65 & 39,3 & 9,3 & Elendi: arşiv radyal hız verisi yetersiz\\\hline
TrES-1 b & 6,3 & 34 & 2,2 & 4,9 & \textbf{Seçildi}\\\hline
TrES-5 b & 4,6 & 66 & 29,2 & 3,2 & Elendi: arşiv radyal hız verisi yetersiz\\\hline
WASP-4 b & 11,5 & 47 & 0,4 & 2,1 & \textbf{Seçildi}\\\hline
WASP-10 b & 7,9 & 8 & – & – & Elendi: TESS geçiş sayısı 30'dan az\\\hline
WASP-12 b & 32,8 & 81 & 1,6 & 10,1 & \textbf{Seçildi}\\\hline
WASP-22 b & 4,2 & 5 & – & – & Elendi: TESS geçiş sayısı 30'dan az\\\hline
WASP-45 b & 8,0 & 14 & – & – & Elendi: TESS geçiş sayısı 30'dan az\\\hline
XO-4 b & 3,6 & 16 & – & – & Elendi: TESS geçiş sayısı 30'dan az\\\hline
HAT-P-7 b & 4,5 & 66 & 1,2 & 6,1 & \textbf{Seçildi}\\\hline
WASP-1 b & 4,3 & 9 & – & – & Elendi: TESS geçiş sayısı 30'dan az\\\hline
WASP-46 b & 18,1 & 32 & 56,3 & 17,2 & Elendi: arşiv radyal hız verisi yetersiz\\\hline
\end{longtable}}
\subsection*{6. Ekler — EK-1: Kaynaklar}
\begin{itemize}[leftmargin=1em,itemsep=1pt]
\item Adams, E. R. vd. 2024. "Doomed Worlds. I. No New Evidence for Orbital Decay in a Long-term Survey of 43 Ultrahot Jupiters", The Planetary Science Journal, 5, 163. \url{https://doi.org/10.3847/PSJ/ad3e80}
\item Akeson, R. L. vd. 2013. "The NASA Exoplanet Archive: Data and Tools for Exoplanet Research", Publications of the Astronomical Society of the Pacific, 125, 989-999. \url{https://doi.org/10.1086/672273}
\item Applegate, J. H. 1992. "A Mechanism for Orbital Period Modulation in Close Binaries", The Astrophysical Journal, 385, 621. \url{https://doi.org/10.1086/170967}
\item Astropy Collaboration vd. 2022. "The Astropy Project: Sustaining and Growing a Community-oriented Open-source Project and the Latest Major Release (v5.0) of the Core Package", The Astrophysical Journal, 935, 167. \url{https://doi.org/10.3847/1538-4357/ac7c74}
\item Baştürk, Ö. vd. 2025. "The Orbit of WASP-4 b is in Decay", Monthly Notices of the Royal Astronomical Society, 541, 714-730. \url{https://doi.org/10.1093/mnras/staf1009}
\item Bonomo, A. S. vd. 2017. "The GAPS Programme with HARPS-N at TNG. XIV. Investigating Giant Planet Migration History via Improved Eccentricity and Mass Determination for 231 Transiting Planets", Astronomy \& Astrophysics, 602, A107. \url{https://doi.org/10.1051/0004-6361/201629882}
\item Bouma, L. G. vd. 2020. "WASP-4 Is Accelerating toward the Earth", The Astrophysical Journal Letters, 893, L29. \url{https://doi.org/10.3847/2041-8213/ab8563}
\item DACE (Data \& Analysis Center for Exoplanets), Cenevre Üniversitesi. \url{https://dace.unige.ch,} Son erişim tarihi: 7 Kasım 2026.
\item Foreman-Mackey, D. vd. 2013. "emcee: The MCMC Hammer", Publications of the Astronomical Society of the Pacific, 125, 306-312. \url{https://doi.org/10.1086/670067}
\item Hagey, S. R. vd. 2025. "TrES-1 b: A Case Study in Detecting Secular Evolution of Exoplanet Orbits", The Astronomical Journal, 170, 197. \url{https://doi.org/10.3847/1538-3881/aded15}
\item Ivshina, E. S. ve Winn, J. N. 2022. "TESS Transit Timing of Hundreds of Hot Jupiters", The Astrophysical Journal Supplement Series, 259, 62. \url{https://doi.org/10.3847/1538-4365/ac545b}
\item Knutson, H. A. vd. 2014. "Friends of Hot Jupiters. I. A Radial Velocity Search for Massive, Long-period Companions to Close-in Gas Giant Planets", The Astrophysical Journal, 785, 126. \url{https://doi.org/10.1088/0004-637X/785/2/126}
\item Kokori, A. vd. 2023. "ExoClock Project. III. 450 New Exoplanet Ephemerides from Ground and Space Observations", The Astrophysical Journal Supplement Series, 265, 4. \url{https://doi.org/10.3847/1538-4365/ac9da4}
\item Kreidberg, L. 2015. "batman: BAsic Transit Model cAlculatioN in Python", Publications of the Astronomical Society of the Pacific, 127, 1161-1165. \url{https://doi.org/10.1086/683602}
\item Kutluay, A. C. vd. 2026. "Chasing the Tides: Searching for Orbital Decay Signatures in Transit Timing Data and Tidal Models for 20 Hot Jupiters", Publications of the Astronomical Society of the Pacific, 138, 024405. \url{https://doi.org/10.1088/1538-3873/ae406a}
\item Lomb, N. R. 1976. "Least-squares Frequency Analysis of Unequally Spaced Data", Astrophysics and Space Science, 39, 447-462. \url{https://doi.org/10.1007/BF00648343}
\item MAST (Mikulski Archive for Space Telescopes), Space Telescope Science Institute. \url{https://archive.stsci.edu,} Son erişim tarihi: 7 Kasım 2026.
\item Oberst, T. E. vd. 2017. "KELT-16b: A Highly Irradiated, Ultra-short Period Hot Jupiter Nearing Tidal Disruption", The Astronomical Journal, 153, 97. \url{https://doi.org/10.3847/1538-3881/153/3/97}
\item Ochsenbein, F., Bauer, P. ve Marcout, J. 2000. "The VizieR Database of Astronomical Catalogues", Astronomy and Astrophysics Supplement Series, 143, 23-32. \url{https://doi.org/10.1051/aas:2000169}
\item Patra, K. C. vd. 2017. "The Apparently Decaying Orbit of WASP-12b", The Astronomical Journal, 154, 4. \url{https://doi.org/10.3847/1538-3881/aa6d75}
\item Penev, K. vd. 2018. "Empirical Tidal Dissipation in Exoplanet Hosts from Tidal Spin-up", The Astronomical Journal, 155, 165. \url{https://doi.org/10.3847/1538-3881/aaaf71}
\item Poddaný, S., Brát, L. ve Pejcha, O. 2010. "Exoplanet Transit Database. Reduction and Processing of the Photometric Data of Exoplanet Transits", New Astronomy, 15, 297-301. \url{https://doi.org/10.1016/j.newast.2009.09.001}
\item Ricker, G. R. vd. 2015. "Transiting Exoplanet Survey Satellite", Journal of Astronomical Telescopes, Instruments, and Systems, 1, 014003. \url{https://doi.org/10.1117/1.JATIS.1.1.014003}
\item Rothmeier, M. vd. 2025. "Doomed Worlds. II. Reassessing Suggestions of Orbital Decay for TrES-5 b", The Planetary Science Journal, 6, 292. \url{https://doi.org/10.3847/PSJ/ae1b9c}
\item Scargle, J. D. 1982. "Studies in Astronomical Time Series Analysis. II. Statistical Aspects of Spectral Analysis of Unevenly Spaced Data", The Astrophysical Journal, 263, 835. \url{https://doi.org/10.1086/160554}
\item Schwarz, G. 1978. "Estimating the Dimension of a Model", The Annals of Statistics, 6, 461-464. \url{https://doi.org/10.1214/aos/1176344136}
\item Sodickson, N. ve Grunblatt, S. 2025. "In Search of Decay: An Analysis of Transit Times of Hot Jupiters in Main-sequence and Post-main-sequence Systems", The Astrophysical Journal, 993, 78. \url{https://doi.org/10.3847/1538-4357/adfd60}
\item T.C. Cumhurbaşkanlığı Strateji ve Bütçe Başkanlığı. 2023. "On İkinci Kalkınma Planı (2024-2028)". \url{https://www.sbb.gov.tr/wp-content/uploads/2023/12/On-Ikinci-Kalkinma-Plani_2024-2028_11122023.pdf,} Son erişim tarihi: 7 Kasım 2026.
\item T.C. Sanayi ve Teknoloji Bakanlığı. 2025. "2030 Sanayi ve Teknoloji Stratejisi". \url{https://www.sanayi.gov.tr/assets/img/popup-image/pdf/SanayiTeknolojiStratejisi2030.pdf,} Son erişim tarihi: 7 Kasım 2026.
\item Wang, W. vd. 2024. "Long-term Variations in the Orbital Period of Hot Jupiters from Transit-timing Analysis Using TESS Survey Data", The Astrophysical Journal Supplement Series, 270, 14. \url{https://doi.org/10.3847/1538-4365/ad0847}
\item Watson, C. A. ve Marsh, T. R. 2010. "Orbital Period Variations of Hot Jupiters Caused by the Applegate Effect", Monthly Notices of the Royal Astronomical Society, 405, 2037. \url{https://doi.org/10.1111/j.1365-2966.2010.16602.x}
\item Winn, J. N. vd. 2009. "HAT-P-7: A Retrograde or Polar Orbit, and a Third Body", The Astrophysical Journal, 703, L99-L103. \url{https://doi.org/10.1088/0004-637X/703/2/L99}
\item Winn, J. N. ve Stefánsson, G. 2025. "Orbital Decay Candidates Reconsidered: WASP-4 b Is Not Decaying and Kepler-1658 b Is Not a Planet", The Planetary Science Journal, 6, 300. \url{https://doi.org/10.3847/PSJ/ae21db}
\item Yang, E., Su, Y. ve Winn, J. N. 2025. "A Third Star in the HAT-P-7 System and a New Dynamical Pathway to Misaligned Hot Jupiters", The Astrophysical Journal, 986, 117. \url{https://doi.org/10.3847/1538-4357/add5f7}
\item Yee, S. W. vd. 2020. "The Orbit of WASP-12b Is Decaying", The Astrophysical Journal Letters, 888, L5. \url{https://doi.org/10.3847/2041-8213/ab5c16}
\end{itemize}


\end{document}
